\chapter{集合论}
\label{cha:set-math}

\index{set|(}%

我们将集合视为具有特别简单同伦性质的类型，这一观念（参见\cref{sec:basics-sets}）与 Zermelo--Fraenkel\index{set theory!Zermelo--Fraenkel} 集合论中的集合颇为不同，后者构成一个具有复杂嵌套隶属结构的累积层次。
对许多数学目的而言，同伦类型论中的集合与 Zermelo--Fraenkel 集合论中的集合同样适用，但二者之间也存在重要的差异。

本章我们首先在\cref{sec:piw-pretopos}中证明，范畴 $\uset$ 拥有（大多数）集合范畴通常具有的性质。
\index{mathematics!constructive}%
\index{mathematics!predicative}%
在构造性、直谓、泛等的基础中，它是一个“$\Pi\mathsf{W}$-预意象（$\Pi\mathsf{W}$-pretopos）”；
\index{propositional!resizing}%
而如果我们假设命题大小重整（\cref{subsec:prop-subsets}），它就是一个初等意象\index{topos}；
如果我们再假设 \LEM{} 和 \choice{}，那么它便是 Lawvere（劳维尔）的\emph{集合范畴的初等理论}\index{Lawvere} 的一个模型。
\index{Elementary Theory of the Category of Sets}%
这足以保证同伦类型论中的集合，其行为与大多数集合论以外的数学家所使用的集合一致。

在本章余下部分，我们将考察一些传统上归入“集合论”的课题。
在\cref{sec:cardinals,sec:ordinals,sec:wellorderings}中，我们将研究基数与序数。
这些概念在集合论中传统上借助全局隶属关系来定义，但我们将会看到，泛等公理能够提供一种同样便捷、更具“结构性”的方法。

最后，在\cref{sec:cumulative-hierarchy}中，我们将探讨在同伦类型论\emph{内部}构造一个累积层次的可能性，并为其配备一个类似于 Zermelo--Fraenkel 集合论的二元隶属关系。
这将高阶归纳类型与代数集合论的思想结合起来。
\index{algebraic set theory}%
\index{set theory!algebraic}%

在本章中，我们将频繁使用\cref{subsec:prop-trunc}中描述的传统逻辑记号。
除了\cref{cha:basics,cha:logic}的基础理论外，我们还会用到\cref{sec:colimits,sec:set-quotients}中关于余极限和商的高阶归纳类型，以及\cref{cha:hlevels}中的一些截断理论，特别是当 $n=-1$ 时\cref{sec:image-factorization}中的因子分解系统。
在\cref{sec:ordinals}中，我们将用归纳族（\cref{sec:generalizations}）来描述良基性；
而在\cref{sec:cumulative-hierarchy}中，则会使用一个更复杂的高阶归纳类型来呈现累积层次。

%\section{\texorpdfstring{$\set$}{Set} is a \texorpdfstring{$\Pi$}{Π}W-pretopos}
\section{The category of sets}
\label{sec:piw-pretopos}

回顾在\cref{cha:category-theory}中，我们将范畴 \uset 定义为由所有 $0$-类型（位于某个宇宙 \UU 中）及其间的映射构成，并指出它是一个范畴（而不仅仅是预范畴）。
下面我们依次考察 \uset 所具备的各级结构。

\subsection{Limits and colimits}
\label{subsec:limits-sets}

\index{limit!of sets}%
\index{colimit!of sets}%

由于集合对积封闭，\cref{thm:prod-ump}中积的泛性质立即表明 \uset 具有有限积。
事实上，由等价 \[ \Parens{X\to \prd{a:A} B(a)} \eqvsym \Parens{\prd{a:A} (X\to B(a))}.\] 同样容易得到无限积。
并且我们在\cref{ex:pullback}\index{pullback}中看到，$f:A\to C$ 和 $g:B\to C$ 的拉回可以定义为 $\sm{a:A}{b:B} f(a)=g(b)$；当 $A,B,C$ 是集合时，它也是集合，且继承了应有的泛性质。
因此，\uset 在自然的意义上是一个\emph{完备的}范畴。
\index{category!complete}%
\index{complete!category}%

由于集合对 $+$ 封闭且包含 \emptyt，\uset 具有有限余积。
类似地，由于当 $A$ 和每个 $B(a)$ 是集合时，$\sm{a:A}B(a)$ 也是集合，所以它在 \uset 中提供了族 $B$ 的余积。
最后，我们在\cref{sec:pushouts}中证明了 $n$-类型中存在推出，特别地，这包含了 \uset。
因此，\uset 也是\emph{余完备的}。
\index{category!cocomplete}%
\index{cocomplete category}%

\subsection{Images}
\label{sec:image}

%We will show that $\uset$ is a $\Pi$W-pretopos.
接下来，我们证明 $\uset$ 是一个\define{正则范畴（regular category）}，即：
\indexdef{category!regular}%
\indexdef{regular!category}%
%

\begin{enumerate}
\item $\uset$ 是有限完备的。\label{item:reg1}
\item 任何函数 $f : A \to B$ 的核偶对 $\proj1,\proj2: (\sm{x,y:A} f(x)= f(y)) \to A$ 有一个余等子。\label{item:reg2}
  \indexdef{kernel!pair}
\item 正则满态射的拉回仍是正则满态射。\label{item:reg3}
\end{enumerate}

%
回忆一下，一个\define{正则满态射（regular epimorphism）}
\indexdef{epimorphism!regular}%
\indexdef{regular!epimorphism}%
是指可以成为\emph{某个}映射对的余等子的态射。
因此，在~\ref{item:reg3}中，要求余等子的拉回仍然是余等子，但不一定是那个被拉回的映射对的余等子。

\index{set-coequalizer}%
\index{image}
对于函数 $f:A\to B$ 的核偶对的余等子，一个自然的候选是 $f$ 的\emph{像（image）}，正如在 \cref{sec:image-factorization} 中所定义的那样。
回顾，我们曾定义 $\im(f)\defeq \sm{b:B} \brck{\hfib f b}$，配有函数
$\tilde{f}:A\to\im(f)$ 与 $i_f:\im(f)\to B$，由
\begin{align*}
  \tilde{f} & \defeq \lam{a} \Pairr{f(a),\,\bproj{\pairr{a,\refl{f(a)}}}}\\
i_f & \defeq \proj1
\end{align*}
定义，构成如下图表：
\begin{equation*}
  \xymatrix{
    **[l]{\sm{x,y:A} f(x)= f(y)}
    \ar@<0.25em>[r]^{\proj1}
    \ar@<-0.25em>[r]_{\proj2}
    &
    {A}
    \ar[r]^(0.4){\tilde{f}}
    \ar[rd]_{f}
    &
    {\im(f)}
    \ar@{..>}[d]^{i_f}
    \\ & &
    B
  }
\end{equation*}

回忆一下，一个函数 $f:A\to B$ 被称为\emph{满的（surjective）}，如果
\index{function!surjective}%
\narrowequation{\fall{b:B}\brck{\hfib f b},}
或等价地 $\fall{b:B} \exis{a:A} f(a)=b$。
我们也曾提到，集合之间的一个函数 $f:A\to B$ 被称为\emph{单的（injective）}，如果
\index{function!injective}%
$\fall{a,a':A} (f(a) = f(a')) \Rightarrow (a=a')$，或等价地，如果它的每个纤维都是命题。
由于这些分别是 \cref{cha:hlevels} 意义下的 $(-1)$-连通映射与 $(-1)$-截断映射，那里的一般理论表明：上图中的 $\tilde f$ 是满的而 $i_f$ 是单的，并且这个分解在拉回下是稳定的。

现在我们将满射性与单射性同相应的范畴论概念等同起来。
首先我们观察到，范畴论中的单态射与满态射有一个稍强的等价表述。

\begin{lem}\label{thm:mono}
  对于范畴 $A$ 中的一个态射 $f:\hom_A(a,b)$，以下条件等价。
  \begin{enumerate}
  \item $f$ 是一个\define{单态射（monomorphism）}：
    \indexdef{monomorphism}%
    对所有 $x:A$ 与 ${g,h:\hom_A(x,a)}$，若 $f\circ g = f\circ h$ 则 $g=h$。\label{item:mono1}
  \item （若 $A$ 有拉回）对角态射 $a\to a\times_b a$ 是一个同构。\label{item:mono4}
  \item 对所有 $x:A$ 和 $k:\hom_A(x,b)$，类型 $\sm{h:\hom_A(x,a)} (k = f\circ h)$ 是一个命题。\label{item:mono2}
  \item 对所有 $x:A$ 与 ${g:\hom_A(x,a)}$，类型 $\sm{h:\hom_A(x,a)} (f\circ g = f\circ h)$ 是可缩的。\label{item:mono3}
  \end{enumerate}
\end{lem}

\begin{proof}
  条件~\ref{item:mono1} 与条件~\ref{item:mono4} 的等价性是标准的范畴论结果。
  现在考虑集合之间的函数 $(f\circ \blank ):\hom_A(x,a) \to \hom_A(x,b)$。
  条件~\ref{item:mono1} 说它是单的，而条件~\ref{item:mono2} 说它的纤维是命题；因此它们等价。
  条件~\ref{item:mono2} 蕴涵条件~\ref{item:mono3}：取 $k\defeq f\circ g$，并注意一个有实例的命题是可缩的。
  最后，由条件~\ref{item:mono3} 可推出条件~\ref{item:mono1}，因为如果 $p:f\circ g= f\circ h$，那么 $(g,\refl{})$ 和 $(h,p)$ 都是条件~\ref{item:mono3} 中类型的项，因此它们相等，于是 $g=h$。
\end{proof}

\begin{lem}\label{thm:inj-mono}
  集合之间的一个函数 $f:A\to B$ 是单的，当且仅当它是 \uset 中的一个单态射\index{monomorphism}。
\end{lem}

\begin{proof}
  留给读者。
\end{proof}

当然，一个\define{满态射（epimorphism）}
\indexdef{epimorphism}%
\indexsee{epi}{epimorphism}%
就是反范畴中的单态射。
我们现在证明，在 \uset 中，满态射恰好就是满射，并且也恰好是余等子（正则满态射）。

$\uset$ 中一对映射 $f,g:A\to B$ 的余等子被定义为一般（同伦）余等子的 0-截断。
为清晰起见，我们可称之为\define{集合余等子（set-coequalizer）}。
\indexdef{set-coequalizer}%
\indexsee{coequalizer!of sets}{set-coequalizer}%
其泛性质可方便地表述如下。

\begin{lem}
\index{universal!property!of set-coequalizer}%
设 $f,g:A\to B$ 是集合 $A$ 与 $B$ 之间的函数。其集合余等子 $c_{f,g}:B\to Q$ 具有以下性质：对任意集合 $C$ 以及任意满足 $h\circ f = h\circ g$ 的 $h:B\to C$，类型
\begin{equation*}
\sm{k:Q\to C} (k\circ c_{f,g} = h)
\end{equation*}
是可缩的。
\end{lem}

\begin{lem}\label{epis-surj}
对于集合之间的任意函数 $f:A\to B$，以下条件等价：
\begin{enumerate}
\item $f$ 是一个满态射。
\item 考虑 $\uset$ 中定义映射锥\index{cone!of a function}的推出图表
\begin{equation*}
  \xymatrix{
    {A}
    \ar[r]^{f}
    \ar[d]
    &
    {B}
    \ar[d]^{\iota}
    \\
    {\unit}
    \ar[r]_{t}
    &
    {C_f}
  }
\end{equation*}
则类型 $C_f$ 是可缩的。
\item $f$ 是满射。
\end{enumerate}
\end{lem}

\begin{proof}
设 $f:A\to B$ 为集合之间的函数，并假设它是满态射；我们证明 $C_f$ 是可缩的。
$C_f$ 的构造子 $\unit\to C_f$ 给出了一个元素 $t:C_f$。
我们需要证明
\begin{equation*}
\prd{x:C_f} x= t.
\end{equation*}
注意 $x= t$ 是一个命题，因此我们可以对 $C_f$ 使用归纳法。
当然，当 $x$ 就是 $t$ 时，我们有 $\refl{t}:t=t$，所以只需找到
\begin{align*}
I_0 & : \prd{b:B} \iota(b)= t\\
I_1 & : \prd{a:A} \opp{\alpha_1(a)} \ct I_0(f(a))=\refl{t}
\end{align*}
其中 $\iota:B\to C_f$ 和 $\alpha_1:\prd{a:A} \iota(f(a))= t$ 是 $C_f$ 的其他构造子。
注意 $\alpha_1$ 是一个从 $\iota\circ f$ 到 $\mathsf{const}_t\circ f$ 的同伦，因此我们得到元素
\begin{equation*}
\pairr{\iota,\refl{\iota\circ f}},\pairr{\mathsf{const}_t,\alpha_1}:
\sm{h:B\to C_f} \iota\circ f \htpy h\circ f.
\end{equation*}
根据 \cref{thm:mono}\ref{item:mono3} 的对偶（以及函数外延性），存在一条道路
\begin{equation*}
\gamma:\pairr{\iota,\refl{\iota\circ f}}=\pairr{\mathsf{const}_t,\alpha_1}.
\end{equation*}
因此，我们可以定义 $I_0(b)\defeq \happly(\projpath1(\gamma),b):\iota(b)=t$。
我们还有
\[\projpath2(\gamma) : \trans{\projpath1(\gamma)}{\refl{\iota\circ f}} = \alpha_1. \]
这个传送涉及与 $f$ 的预复合，而预复合与 $\happly$ 交换。
因此，由道路类型中的传送，对任意 $a:A$ 我们得到 $I_0(f(a)) = \alpha_1(a)$，这就给出了 $I_1$。

现在假设 $C_f$ 是可缩的；我们证明 $f$ 是满的。
我们首先通过对 $C_f$ 递归定义一个类型族 $P:C_f\to\prop$，这是可行的，因为 \prop 是一个集合。
在点构造子上，我们定义
\begin{align*}
P(t) & \defeq \unit\\
P(\iota(b)) & \defeq \brck{\hfiber{f}b}.
\end{align*}
为了完成 $P$ 的构造，还需要给出一条道路
\narrowequation{\brck{\hfiber{f}{f(a)}} =_\prop \unit}
对所有 $a:A$。
然而，$\brck{\hfiber{f}{f(a)}}$ 有实例 $(f(a),\refl{f(a)})$。
由于它是一个命题，这意味着它是可缩的——从而等价，因此相等，于 \unit。
这就完成了 $P$ 的定义。
现在，既然 $C_f$ 被假设为可缩的，那么对任意 $x:C_f$，$P(x)$ 等价于 $P(t)$。
特别地，对每个 $b:B$，$P(\iota(b))\jdeq \brck{\hfiber{f}b}$ 等价于 $P(t)\jdeq \unit$，因此是可缩的。
于是，$f$ 是满的。

最后，假设 $f:A\to B$ 是满的，并考虑一个集合 $C$ 以及两个函数 $g,h:B\to C$ 满足 $g\circ f = h\circ f$。
由于 $f$ 被假设为满的，对所有 $b:B$，类型 $\brck{\hfib f b}$ 是可缩的。
因此我们有如下等价：
\begin{align*}
\prd{b:B} (g(b)= h(b))
& \eqvsym \prd{b:B} \Parens{\brck{\hfib f b} \to (g(b)= h(b))}\\
& \eqvsym \prd{b:B} \Parens{\hfib f b \to (g(b)= h(b))}\\
& \eqvsym \prd{b:B}{a:A}{p:f(a)= b} g(b)= h(b)\\
& \eqvsym \prd{a:A} g(f(a))= h(f(a))
\end{align*}
在第二行中我们用到了 $g(b)=h(b)$ 是一个命题这一事实，因为 $C$ 是一个集合。
但由假设，后一个类型是有实例的。
\end{proof}

% \begin{rem}
% The above theorem is not true when we replace $\set$ by $\type$
% (replacing it also in the definition of $\mathsf{epi}$ and $\mathsf{epi}'$).
% However, we do
% get the implications $\textit{ii.}\Rightarrow\textit{iii.}\Rightarrow
% \textit{iv.}$
% \end{rem}

\begin{thm}\label{thm:set_regular}\label{lem:images_are_coequalizers}
范畴 $\uset$ 是正则范畴。此外，集合之间的满函数是正则满态射。
\end{thm}

\begin{proof}
范畴论中有一个标准引理：如果一个范畴有有限极限，并且有一个拉回稳定的正交分解系统\index{orthogonal factorization system} $(\mathcal{E},\mathcal{M})$，其中 $\mathcal{M}$ 是单态射，那么该范畴是正则的，此时 $\mathcal{E}$ 自动由正则满态射组成。
（例如参见~\cite[A1.3.4]{elephant}。）
该分解系统的存在性已在 \cref{thm:orth-fact} 中证明。
\end{proof}

\begin{lem}\label{lem:pb_of_coeq_is_coeq}
在 \uset 中，正则满态射的拉回仍是正则满态射。
\end{lem}

\begin{proof}
  我们在 \cref{thm:stable-images} 中证明了 $n$-连通函数的拉回是 $n$-连通的。
  根据 \cref{lem:images_are_coequalizers}，只需取 $n=-1$ 应用此结论即可。
\end{proof}

\indexdef{image!of a subset}
作为 \uset（集合的范畴）是正则范畴的一个推论，我们可以在子集上定义“像”操作。
也就是说，给定 $f:A\to B$，$A$ 的任意子集 $P:\power A$（即一个谓词 $P:A\to \prop$）都有一个\define{像（image）}，它是 $B$ 的一个子集。
这可以直接定义为 $\setof{ y:B | \exis{x:A} f(x)=y \land P(x)}$，或者间接地定义为复合函数
\[ \setof{ x:A | P(x) } \to A \xrightarrow{f} B.\]
\symlabel{subset-image}
的像（按前述意义）。
我们有时也会使用常见的记号 $\setof{f(x) | P(x)}$ 来表示 $P$ 的像。

\subsection{商}\label{subsec:quotients}

\index{set-quotient|(}%
既然我们已经知道 $\uset$ 是正则的，那么为了证明 $\uset$ 是正合的，就需要证明每个等价关系都是有效的。
\index{effective!equivalence relation|(}%
\index{relation!effective equivalence|(}%
换句话说，给定一个等价关系 $R:A\to A\to\prop$，存在一对映射
$\proj1,\proj2:\sm{x,y:A} R(x,y)\to A$
的余等子 $c_R$；并且，$\proj1$ 和 $\proj2$ 构成了 $c_R$ 的核偶\index{kernel!pair}。

我们已经在 \cref{sec:set-quotients} 中见过两种由等价关系 $R:A\to A\to\prop$ 构造集合商的一般方法。
第一种可以描述为两个投影
\[\proj1,\proj2:\Parens{\sm{x,y:A} R(x,y)} \to A.\]
的集合余等子。
这类商的一个重要性质如下所述。

\begin{defn}
  如果下面的方块
\begin{equation*}
  \xymatrix{
    {\sm{x,y:A} R (x,y)}
    \ar[r]^(0.7){\proj1}
    \ar[d]_{\proj2}
    &
    {A}
    \ar[d]^{c_R}
    \\
    {A}
    \ar[r]_{c_R}
    &
    {A/R}
    }
\end{equation*}
是一个拉回，那么关系 $R:A\to A\to\prop$ 就被称为是\define{有效的}。
\indexdef{effective!relation}
\indexdef{effective!equivalence relation}%
\indexdef{relation!effective equivalence}%
\end{defn}

由于 $c_R$ 和自身的标准拉回是 $\sm{x,y:A} (c_R(x)=c_R(y))$，根据 \cref{thm:total-fiber-equiv}，这等价于要求典范变换 $\prd{x,y:A} R(x,y) \to (c_R(x)=c_R(y))$ 是一个逐纤维等价。

\begin{lem}\label{lem:sets_exact}
假设 $\pairr{A,R}$ 是一个等价关系。那么对任意 $x,y:A$，存在一个等价
\begin{equation*}
(c_R(x)= c_R(y))\eqvsym R(x,y)
\end{equation*}
换言之，等价关系是有效的。
\end{lem}

\begin{proof}
我们首先将 $R$ 延拓成一个关系 $\widetilde{R}:A/R\to A/R\to\prop$，随后我们将证明它等价于 $A/R$ 上的相等类型。我们通过对 $A/R$ 进行双重归纳来定义 $\widetilde{R}$（注意，根据命题的泛等公理，$\prop$ 是一个集合）。我们定义 $\widetilde{R}(c_R(x),c_R(y)) \defeq R(x,y)$。对于 $r:R(x,x')$ 和 $s:R(y,y')$，利用 $R$ 的传递性与对称性，我们可以从 $R(x,y)$ 到 $R(x',y')$ 给出一个等价。这就完成了 $\widetilde{R}$ 的定义。

剩下只需证明，对每一对 $w,w':A/R$，都有 $\widetilde{R}(w,w')\eqvsym (w= w')$。
方向 $(w=w')\to \widetilde{R}(w,w')$ 可以通过传送得到，只要我们证明 $\widetilde{R}$ 是自反的，而这可以通过一个简单的归纳完成。
另一个方向 $\widetilde{R}(w,w')\to (w= w')$ 是一个命题，并且由于 $c_R:A\to A/R$ 是满的，我们只需假设 $w$ 和 $w'$ 具有 $c_R(x)$ 和 $c_R(y)$ 的形式。
但在这种情况下，我们已有典范映射 $\widetilde{R}(c_R(x),c_R(y)) \defeq R(x,y) \to (c_R(x)=c_R(y))$。
（再次注意，这里出现了编码-解码方法。\index{encode-decode method}）
\end{proof}

商的第二种构造是作为 $R$ 的等价类组成的集合（即其幂集\index{power set}的一个子集）：
\[ A\sslash R \defeq \setof{ P:A\to\prop | P \text{ is an equivalence class of } R}. \]
为了使该构造与 $A$ 和 $R$ 保持在同一个宇宙中，需要命题大小重整\index{propositional resizing}\index{impredicative!quotient}\index{resizing}。

注意，如果我们把 $R$ 看作从 $A$ 到 $A\to \prop$ 的函数，那么 $A\sslash R$ 就等价于 \cref{sec:image} 中构造的像 $\im(R)$。
现在，在 \cref{lem:images_are_coequalizers} 中我们已经证明了像是余等子。特别地，我们立刻得到余等子图表
\begin{equation*}
  \xymatrix{
    **[l]{\sm{x,y:A} R (x)= R (y)}
    \ar@<0.25em>[r]^{\proj1}
    \ar@<-0.25em>[r]_{\proj2}
    &
    {A}
    \ar[r]
    &
    {A \sslash R.}
  }
\end{equation*}
我们可以利用这个来给出另一个证明：任何等价关系都是有效的，并且两种商的定义是一致的。

\begin{thm}\label{prop:kernels_are_effective}
对于任意两个集合之间的任意函数 $f:A\to B$，由
$\ker(f,x,y)\defeq (f(x)= f(y))$
给出的关系 $\ker(f):A\to A\to\prop$ 是有效的。
\end{thm}

\begin{proof}
我们将使用 $\im(f)$ 是 $\proj1,\proj2:
(\sm{x,y:A} f(x)= f(y))\to A$ 的余等子这一事实。
%we get this equivalence from~\cref{prop:images_are_coequalizers}
注意，函数
\[c_f\defeq\lam{a} \Parens{f(a),\brck{\pairr{a,\refl{f(a)}}}}
: A \to \im(f)
\]
的核偶由两个投影
\begin{equation*}
\proj1,\proj2:\Parens{\sm{x,y:A} c_f(x)= c_f(y)}\to A.
\end{equation*}
组成。
对于任意 $x,y:A$，我们有等价
\begin{align*}
  (c_f(x)= c_f(y))
  & \eqvsym \Parens{\sm{p:f(x)= f(y)} \trans{p}{\brck{\pairr{x,\refl{f(x)}}}} =\brck{\pairr{y,\refl{f(x)}}}}\\
  & \eqvsym (f(x)= f(y)),
\end{align*}
其中最后一个等价成立，是因为对于任意 $b:B$，$\brck{\hfiber{f}b}$ 是一个命题。
因此，我们得到
\begin{equation*}
\Parens{\sm{x,y:A} c_f(x)= c_f(y)}\eqvsym \Parens{\sm{x,y:A} f(x)= f(y)}
\end{equation*}
由此我们可以得出结论：对于任意函数 $f$，其核 $\ker f$ 都是一个有效关系。
\end{proof}

\begin{thm}
等价关系是有效的，并且存在一个等价 $A/R \eqvsym A\sslash  R $。
\end{thm}

\begin{proof}
我们需要分析余等子图表
\begin{equation*}
  \xymatrix{
    **[l]{\sm{x,y:A} R (x)= R (y)}
    \ar@<0.25em>[r]^{\proj1}
    \ar@<-0.25em>[r]_{\proj2}
    &
    {A}
    \ar[r]
    &
    {A \sslash R}
  }
\end{equation*}
根据泛等公理，类型 $R(x) = R(y)$ 等价于从 $R(x)$ 到 $R(y)$ 的同伦的类型，而这个类型又等价于
\narrowequation{\prd{z:A} R (x,z)\eqvsym R (y,z).}
由于 $R$ 是一个等价关系，后一个类型等价于 $R(x,y)$。综上所述，我们得到 $(R(x) = R(y)) \eqvsym R(x,y)$，因此 $R$ 是有效的，因为它等价于一个有效关系。并且，图表
\begin{equation*}
  \xymatrix{
    **[l]{\sm{x,y:A} R(x, y)}
    \ar@<0.25em>[r]^{\proj1}
    \ar@<-0.25em>[r]_{\proj2}
    &
    {A}
    \ar[r]
    &
    {A \sslash R.}
  }
\end{equation*}
是一个余等子图表。由于余等子在等价意义下是唯一的，于是有 $A/R \eqvsym A\sslash  R $。
\end{proof}

我们以商的第三种可能的构造来结束本节。设 $A$ 是一个集合，$R$ 是其上的一个等价关系。考虑以 $A$ 为对象、以 $R$ 为态射集的预范畴；那么该预范畴的 Rezk 完备化
\index{completion!Rezk}%
（参见 \cref{sec:rezk}）的对象类型就是这个商。请读者自行验证其中的细节。

\index{effective!equivalence relation|)}%
\index{relation!effective equivalence|)}%
\index{set-quotient|)}%

\subsection{\texorpdfstring{$\uset$}{Set} 是一个 \texorpdfstring{$\Pi\mathsf{W}$}{ΠW}-预意象}
\label{subsec:piw}

\index{structural!set theory|(}%

所谓 \emph{$\Pi\mathsf{W}$-预意象}
\index{PiW-pretopos@$\Pi\mathsf{W}$-pretopos}%
\indexsee{pretopos}{$\Pi\mathsf{W}$-预意象}
——即一个具有不相交有限余积、有效等价关系以及多项式自函子之初始代数的局部笛卡尔闭范畴
\index{locally cartesian closed category}%
\index{category!locally cartesian closed}%
——旨在作为一种“直谓的”
\index{mathematics!predicative}%
意象概念，亦即一个“直谓集合”的范畴。它可以为构造性数学
\index{mathematics!constructive}%
提供基础，其作用类似于通常的集合范畴之于经典
\index{mathematics!classical}%
数学。

通常，在构造性类型论中，人们需要诉诸于一种外部的 ``setoid'' 构造——即正合完备化——以得到一个具有此类闭包性质的范畴。
\index{setoid}\index{completion!exact}%
  特别是，数学中许多通常涉及（非构造性）幂集的构造都需要性质良好的商。值得注意的是，泛等基础能够\emph{在内部}（通过高阶归纳类型）提供这些构造，而无须借助此类外部构造。这体现了我们方法的一项显著优势，在后续的例子中我们将看到这一点。

\begin{thm}
  \index{PiW-pretopos@$\Pi\mathsf{W}$-预意象}
  范畴 $\uset$ 是一个 $\Pi\mathsf{W}$-预意象（$\Pi\mathsf{W}$-pretopos）。
\end{thm}

\begin{proof}
  我们有一个初始对象
  \index{initial!set}%
  $\emptyt$ 以及有限的、不相交的和 $A+B$。  这些结构在拉回下是稳定的，原因仅在于拉回具有右伴随\index{adjoint!functor}。  事实上，$\uset$ 是局部笛卡尔闭的：因为对于集合间的任意映射 $f:A\to B$，其``纤维化替换''\index{fibrant replacement} $\sm{a:A}f(a)=b$ 在 $B$ 之上等价于 $A$，并且该替换具有依值函数类型。

我们刚刚已经证明了 $\uset$ 是正则的（\cref{thm:set_regular}）并且其商是有效的（\cref{lem:sets_exact}）。因此，我们得到了一个局部笛卡尔闭的预意象。最后，根据\cref{ex:ntypes-closed-under-wtypes}，$n$-类型对 $W$-类型的形成是封闭的，并且根据\cref{thm:w-hinit}，$W$-类型是多项式自函子的初始代数。由此可知，$\uset$ 是一个 $\Pi\mathsf{W}$-预意象。
\end{proof}

\index{topos|(}
人们自然会问，是否有某种因素阻碍 $\uset$ 成为（初等）意象？
除了已经提到的结构外，一个意象还具有一个
\emph{子对象分类子（subobject classifier）}：
\indexdef{subobject classifier}%
\index{classifier!subobject}%
\index{power set}%
即一个带点对象，分类单态射（的等价类）。\index{monomorphism}（事实上，有了子对象分类子后，事情会变得更简单一些：只需笛卡尔闭性就能得到余极限。）

在同伦类型论中，泛等性意味着类型$\prop \defeq \sm{X:\UU}\isprop(X)$确实分类单态射（论证类似于\cref{sec:object-classification}），但一般来说，它和背景宇宙 $\UU$ 一样大。
因此，在"是 $0$-类型"的意义上它是一个"集合"，但在"是 $\UU$ 中对象"的意义上它并非"小的"，故而也不属于范畴 \uset。
然而，如果我们假定一种恰当的命题大小重整形式（参见\cref{subsec:prop-subsets}），那么我们就能找到一个 $\prop$ 的小版本，使得 $\uset$ 成为一个初等意象。

\begin{thm}\label{thm:settopos}
  \index{propositional!resizing}%
  如果存在一个包含所有命题的类型 $\Omega:\UU$，那么范畴 $\uset_\UU$ 是一个初等意象。
\end{thm}


\index{topos|)}

为此的一个充分条件是排中律，即我们称之为 \LEM{} 的"命题"形式；因为此时有 $\prop = \bool$，而 $\bool$ \emph{是}小的，并且它也分类所有命题。
此外，在意象理论中，\LEM{} 的一个众所周知的充分条件是选择公理，而选择公理在经典\index{mathematics!classical}集合论中当然通常被假定为公理。
在下一节中，我们将简要探讨在我们的框架下这些条件之间的关系。

\index{structural!set theory|)}%

%%%%%%%%%%%%%%%%%%%%%%%%%%%%%%%%%%%%%%%%%%%%%%%%
\subsection{选择公理蕴含排中律}
\label{subsec:emacinsets}

% In this section we prove a classic result that the axiom of choice implies excluded
% middle.

我们先给出以下引理。

\begin{lem}\label{prop:trunc_of_prop_is_set}
若 $A$ 是一个命题，则其纬悬 $\susp(A)$ 是一个集合，且 $A$ 等价于$\id[\susp(A)]{\north}{\south}$。
\end{lem}

\begin{proof}
为了证明 $\susp(A)$ 是一个集合，我们定义一个族 $P:\susp(A)\to\susp(A)\to\type$，它具有下述性质：对每个 $x,y:\susp(A)$，$P(x,y)$ 是一个命题，并且它等价于恒等类型 $\idtypevar{\susp(A)}$。
%
我们作如下定义：
\begin{align*}
P(\north,\north) & \defeq \unit &
P(\south,\north) & \defeq A\\
P(\north,\south) & \defeq A &
P(\south,\south) & \defeq \unit.
\end{align*}
我们必须验证该定义保持路径。
给定任意 $a : A$，存在一条经线 $\merid(a) : \north = \south$，因此我们也应该有
%
\begin{equation*}
  P(\north, \north) = P(\north, \south) = P(\south, \north) = P(\south, \south).
\end{equation*}
%
但由于 $A$ 有实例 $a$，它等价于 $\unit$，所以我们有
%
\begin{equation*}
  P(\north, \north) \eqvsym P(\north, \south) \eqvsym P(\south, \north) \eqvsym P(\south, \south).
\end{equation*}
%
泛等公理将这些转化为所需的相等性。同时，$P(x,y)$ 对于所有 $x, y : \susp(A)$ 都是一个命题，这可以通过对 $x$ 和 $y$ 进行归纳，并利用"是命题"本身是一个命题这一事实来证明。

注意 $P$ 是一个自反关系。
因此我们可以应用 \cref{thm:h-set-refrel-in-paths-sets}，于是只需构造 $\tau : \prd{x,y:\susp(A)}P(x,y)\to(x=y)$。我们通过双归纳来完成这一构造。
当 $x$ 是 $\north$ 时，我们通过下式定义 $\tau(\north)$：
%
\begin{equation*}
  \tau(\north,\north,u) \defeq \refl{\north}
  \qquad\text{and}\qquad
  \tau(\north,\south,a) \defeq \merid(a).
\end{equation*}
%
如果 $A$ 有实例 $a$，那么 $\merid(a) : \north = \south$，因此我们还需要
\narrowequation{
  \trans{\merid(a)}{\tau(\north, \north)} = \tau(\north, \south).
}
这可由函数外延性得到，利用如下事实：对于所有 $x : A$，
%
\begin{multline*}
  \trans{\merid(a)}{\tau(\north,\north,x)} =
  \tau(\north,\north,x) \ct \opp{\merid(a)} \jdeq \\
  \refl{\north} \ct \merid(a) =
  \merid(a) =
  \merid(x) \jdeq
  \tau(\north, \south, x).
\end{multline*}
对称地，我们可以通过下式定义 $\tau(\south)$：
%
\begin{equation*}
  \tau(\south,\north, a) \defeq \opp{\merid(a)}
  \qquad\text{and}\qquad
  \tau(\south,\south, u) \defeq \refl{\south}.
\end{equation*}
%
为了完成 $\tau$ 的构造，需要验证 $\trans{\merid(a)}{\tau(\north)} = \tau(\south)$，给定任意 $a : A$。验证过程与前类似，通过对 $\tau$ 的第二个参数进行归纳来完成。

因此，由 \cref{thm:h-set-refrel-in-paths-sets}，我们有 $\susp(A)$ 是一个集合，并且对于所有 $x,y:\susp(A)$，$P(x,y) \eqvsym (\id{x}{y})$。
取 $x\defeq \north$ 和 $y\defeq \south$ 即得到所需的 $\eqv{A}{(\id[\susp(A)]\north\south)}$。
\end{proof}

\begin{thm}[Diaconescu（迪亚科内斯库）]\label{thm:1surj_to_surj_to_pem}
  \index{axiom!of choice}%
  \index{excluded middle}%
  \index{Diaconescu's theorem}\index{theorem!Diaconescu's}%
  选择公理蕴含排中律。
\end{thm}

\begin{proof}
我们使用在 \cref{thm:ac-epis-split} 中给出的选择公理的等价形式。
考虑一个命题 $A$。
由 $f(\bfalse) \defeq \north$ 和 $f(\btrue) \defeq \south$ 定义的函数 $f:\bool\to\susp(A)$ 是满的。
事实上，我们有
$\pairr{\bfalse,\refl{\north}} : \hfiber{f}{\north}$
以及 $\pairr{\btrue,\refl{\south}} :\hfiber{f}{\south}$。
由于 $\bbrck{\hfiber{f}{x}}$ 是一个命题，通过归纳即可得证所述的满射性。

由 \cref{prop:trunc_of_prop_is_set}，纬悬 $\susp(A)$ 是一个集合，因此根据选择公理，仅仅存在 $f$ 的一个截面 $g: \susp(A) \to \bool$。
由于 $\bool$ 上的相等性是可判定的，我们得到
\begin{equation*}
 (g(f(\bfalse))= g(f(\btrue))) +
 \lnot (g(f(\bfalse))= g(f(\btrue))),
\end{equation*}
并且，由于 $g$ 是 $f$ 的截面，从而是单的，
\begin{equation*}
(f(\bfalse) = f(\btrue)) +
\lnot (f(\bfalse) = f(\btrue)).
\end{equation*}
最后，因为由 \cref{prop:trunc_of_prop_is_set} 有 $(f(\bfalse)=f(\btrue)) = (\north=\south) = A$，我们得到 $A+\neg A$。
\end{proof}

% This conclusion needs only \LEM{}, see \cref{ex:lemnm}.

% \begin{cor}\label{cor:ACtoLEM0}
%   If the axiom of choice \choice{} holds then $\brck{A + \neg A}$ for every set $A$.
% \end{cor}

% \begin{proof}
%   There is a surjection
%   \[
%   A + \neg A \epi \brck{A} + \brck{\neg A} \epi
%   \brck{(\brck{A} + \brck{\neg A})} = \brck{A} \vee \brck{\neg A} = \brck{A} \vee \neg \brck{A} = \unit,
%   \]
%   %
%   where in the last step excluded middle is available as a consequence of the axiom of choice.
%   Again by the axiom of choice there merely exists a section of the surjection, but this
%   is none other than an inhabitant of $A + \neg A$. Therefore $\brck{A+\neg A}$.
% \end{proof}

\index{denial}


\begin{thm}\label{thm:ETCS}
  \index{Elementary Theory of the Category of Sets}%
  \index{category!well-pointed}%
  如果选择公理成立，那么范畴 $\uset$ 是一个良点化的、带选择公理的布尔初等意象。
\end{thm}

\begin{proof}
  由于 \choice{} 蕴含 \LEM{}，根据 \cref{thm:settopos} 及其后的注记，我们得到一个带选择公理的布尔初等意象。我们将良点性的证明留给读者作为练习（\cref{ex:well-pointed}）。
\end{proof}

\begin{rmk}
  定理中提到的范畴条件被称为 Lawvere（劳维尔）为\emph{集合范畴的初等理论（Elementary Theory of the Category of Sets）}~\cite{lawvere:etcs-long} 所提出的公理。
\end{rmk}

\section{基数}
\label{sec:cardinals}

\begin{defn}\label{defn:card}
  \define{基数类型（type of cardinal numbers）}
  \indexdef{type!of cardinal numbers}%
  \indexdef{cardinal number}%
  \indexsee{number!cardinal}{cardinal number}%
  是集合类型 \set 的 0-截断：
  \[ \card \defeq \pizero{\set} \]
  因此，一个\define{基数（cardinal number）}，或简称\define{基数（cardinal）}，是 $\card\jdeq \pizero\set$ 的一个实例。当然，从技术上讲，每个宇宙 \type 都关联有一个独立的类型 $\card_\UU$。
\end{defn}

%\begin{rmk}

  % , but with these conventions we can state theorems beginning with ``for all cardinal numbers\dots''\ and give them exactly the same sort of meaning as those beginning ``for all types\dots''.
%\end{rmk}

按照截断的惯例，如果 $A$ 是一个集合，则 $\cd{A}$ 表示其在典范投影 $\set \to \trunc0\set \jdeq \card$ 下的像；我们称 $\cd{A}$ 为 $A$ 的\define{基数（cardinality）}\indexdef{cardinality}。
根据定义，\card 是一个集合。
它还从 \set 继承了半环结构。

\begin{defn}
  \define{基数加法（cardinal addition）}
  \indexdef{addition!of cardinal numbers}%
  \index{cardinal number!addition of}%
  \[ (\blank+\blank) : \card \to \card \to \card \]
  通过截断上的归纳来定义：
  \[ \cd{A} + \cd{B} \defeq \cd{A+B}.\]
\end{defn}

\begin{proof}
  由于 $\card\to\card$ 是一个集合，要为所有 $\alpha:\card$ 定义 $(\alpha+\blank):\card\to\card$，由归纳法，只需假设 $\alpha$ 是 $\cd{A}$，其中 $A:\set$。
  现在我们要定义 $(\cd{A}+\blank) :\card\to\card$，即要为所有 $\beta:\card$ 定义 $\cd{A}+\beta :\card$。
  然而，由于 $\card$ 是一个集合，由归纳法只需假设 $\beta$ 是 $\cd{B}$，其中 $B:\set$。
  此时，我们可以将 $\cd{A}+\cd{B}$ 定义为 $\cd{A+B}$。
\end{proof}

\begin{defn}
  类似地，\define{基数乘法}
  \indexdef{multiplication!of cardinal numbers}%
  \index{cardinal number!multiplication of}%
  \[ (\blank\cdot\blank) : \card \to \card \to \card \]
  运算也由截断归纳定义：
  \[ \cd{A} \cdot \cd{B} \defeq \cd{A\times B} \]
\end{defn}

\begin{lem}\label{card:semiring}
  \card 是一个交换半环\index{semiring}，即对任意 $\alpha,\beta,\gamma:\card$，我们有：
  \begin{align*}
    (\alpha+\beta)+\gamma &= \alpha+(\beta+\gamma)\\
    \alpha+0 &= \alpha\\
    \alpha + \beta &= \beta + \alpha\\
    (\alpha \cdot \beta) \cdot \gamma &= \alpha \cdot (\beta\cdot\gamma)\\
    \alpha \cdot 1 &= \alpha\\
    \alpha\cdot\beta &= \beta\cdot\alpha\\
    \alpha\cdot(\beta+\gamma) &= \alpha\cdot\beta + \alpha\cdot\gamma
  \end{align*}
  其中 $0 \defeq \cd{\emptyt}$，$1\defeq\cd{\unit}$。
\end{lem}

\begin{proof}
  我们证明乘法的交换律 $\alpha\cdot\beta = \beta\cdot\alpha$；其他性质的证明完全类似。
  由于 \card 是一个集合，类型 $\alpha\cdot\beta = \beta\cdot\alpha$ 是一个命题，特别地也是一个集合。
  于是，由归纳法，只需假设存在 $A,B:\set$ 使得 $\alpha$ 和 $\beta$ 分别是 $\cd{A}$ 和 $\cd{B}$。
  此时，$\cd{A}\cdot \cd{B} \jdeq \cd{A\times B}$ 且 $\cd{B}\cdot\cd{A} \jdeq \cd{B\times A}$，因此只需证明 $A\times B = B\times A$。
  最后，由泛等性，只需给出一个等价 $A\times B \eqvsym B\times A$。
  这很简单：取 $(a,b) \mapsto (b,a)$ 及其显然的逆即可。
\end{proof}

\begin{defn}
  \define{基数幂}运算同样由截断归纳定义：
  \indexdef{exponentiation, of cardinal numbers}%
  \index{cardinal number!exponentiation of}%
  \[ \cd{A}^{\cd{B}} \defeq \cd{B\to A}. \]
\end{defn}

\begin{lem}\label{card:exp}
  对任意 $\alpha,\beta,\gamma:\card$，有：
  \begin{align*}
    \alpha^0 &= 1\\
    1^\alpha &= 1\\
    \alpha^1 &= \alpha\\
    \alpha^{\beta+\gamma} &= \alpha^\beta \cdot \alpha^\gamma\\
    \alpha^{\beta\cdot \gamma} &= (\alpha^{\beta})^\gamma\\
    (\alpha\cdot\beta)^\gamma &= \alpha^\gamma \cdot \beta^\gamma
  \end{align*}
\end{lem}

\begin{proof}
  证明方法同 \cref{card:semiring}。
\end{proof}

\begin{defn}
  \define{基数不等式}这一关系
  \index{order!non-strict}%
  \index{cardinal number!inequality of}%
  \[ (\blank\le\blank) : \card\to\card\to\prop \]
  由截断归纳定义：
  \symlabel{inj}
  \[ \cd{A} \le \cd{B} \defeq \brck{\inj(A,B)} \]
  其中 $\inj(A,B)$ 是从 $A$ 到 $B$ 的单射的类型。
  \index{function!injective}%
  换句话说，$\cd{A} \le \cd{B}$ 意味着仅知存在从 $A$ 到 $B$ 的单射。
\end{defn}

\begin{lem}
  基数不等式是一个预序，即对于 $\alpha,\beta:\card$，有
  \index{preorder!of cardinal numbers}%
  \begin{gather*}
    \alpha \le \alpha\\
    (\alpha \le \beta) \to (\beta\le\gamma) \to (\alpha\le\gamma)
  \end{gather*}
\end{lem}

\begin{proof}
  与之前一样，由截断归纳。
  以传递性为例：由于 $(\alpha \le \beta) \to (\beta\le\gamma) \to (\alpha\le\gamma)$ 是一个命题，由 0-截断归纳，可以假设 $\alpha$、$\beta$ 和 $\gamma$ 分别是 $\cd{A}$、$\cd{B}$ 和 $\cd{C}$。
  现在，由于 $\cd{A} \le \cd{C}$ 也是一个命题，由 $(-1)$-截断归纳，可以假设给定了单射 $f:A\to B$ 和 $g:B\to C$。
  那么 $g\circ f$ 就是从 $A$ 到 $C$ 的一个单射，因此 $\cd{A} \le \cd{C}$ 成立。
  自反性的证明则更为简单。
\end{proof}

同样地，可以证明基数不等式与半环运算是相容的。

\begin{lem}\label{thm:injsurj}
  \index{function!injective}%
  \index{function!surjective}%
  考虑以下陈述：
  \begin{enumerate}
  \item 存在一个单射 $A\to B$。\label{item:cle-inj}
  \item 存在一个满射 $B\to A$。\label{item:cle-surj}
  \end{enumerate}
  那么，在假设排中律的前提下：
  \index{excluded middle}%
  \index{axiom!of choice}%
  \begin{itemize}
  \item 给定 $a_0:A$，有 \ref{item:cle-inj} $\to$ \ref{item:cle-surj}。
  \item 因此，如果 $A$ 仅知有实例，有 \ref{item:cle-inj} $\to$ 仅知 \ref{item:cle-surj}。
  \item 假设选择公理，有 \ref{item:cle-surj} $\to$ 仅知 \ref{item:cle-inj}。
  \end{itemize}
\end{lem}

\begin{proof}
  如果 $f:A\to B$ 是一个单射，如下定义 $g:B\to A$。
  对于 $b:B$，由于 $f$ 是单的，$f$ 在 $b$ 处的纤维是一个命题。
  因此，根据排中律，要么存在 $a:A$ 使得 $f(a)=b$，要么不存在。
  在第一种情况下，定义 $g(b)\defeq a$；否则令 $g(b)\defeq a_0$。
  那么对于任意 $a:A$，都有 $a = g(f(a))$，所以 $g$ 是满射。

  第二个陈述由第一个陈述经截断归纳得出。
  对于第三个陈述，如果 $g:B\to A$ 是满射，那么根据选择公理，仅知存在一个函数 $f:A\to B$，使得对所有 $a$ 都有 $g(f(a)) = a$。
  而这个 $f$ 必然是单射。
\end{proof}

\begin{thm}[Schroeder--Bernstein]
  \index{theorem!Schroeder--Bernstein}%
  \index{Schroeder--Bernstein theorem}%
  假设排中律，对于集合 $A$ 和 $B$，有
  \[ \inj(A,B) \to \inj(B,A) \to (A\cong B) \]
\end{thm}

\begin{proof}
  通常的"来回"论证可直接适用，无需显著改动。
  注意，它实际上构造了一个同构 $A\cong B$（这里假设了排中律，以便判定一个给定元素是属于一个循环、一条无限链、一条始于 $A$ 的链，还是一条始于 $B$ 的链）。
\end{proof}

\begin{cor}
  假设排中律，基数不等式是一个偏序，即对于 $\alpha,\beta:\card$，有
  \[ (\alpha\le\beta) \to (\beta\le\alpha) \to (\alpha=\beta). \]
\end{cor}

\begin{proof}
  由于 $\alpha=\beta$ 是一个命题，由截断归纳，可以假设 $\alpha$ 和 $\beta$ 分别是 $\cd{A}$ 和 $\cd{B}$，并且有单射 $f:A\to B$ 和 $g:B\to A$。
  那么由 Schroeder--Bernstein 定理，得到一个同构 $A\cong B$，进而得到相等 $\cd{A}=\cd{B}$。
\end{proof}

最后，我们可以复现 Cantor 定理，它表明对于每个基数，都存在一个更大的基数。

\begin{thm}[Cantor]
  \index{Cantor's theorem}%
  \index{theorem!Cantor's}%
  对于 $A:\set$，不存在满射 $A \to (A\to \bool)$。
\end{thm}

\begin{proof}
  假设 $f:A \to (A\to \bool)$ 是任意函数，定义 $g:A\to \bool$ 为 $g(a) \defeq \neg f(a)(a)$。
  如果 $g = f(a_0)$，则 $g(a_0) = f(a_0)(a_0)$ 但 $g(a_0) = \neg f(a_0)(a_0)$，产生矛盾。
  因此，$f$ 不是满射。
\end{proof}

\begin{cor}
  假设排中律，对于任意 $\alpha:\card$，存在基数 $\beta$ 使得 $\alpha\le\beta$ 且 $\alpha\neq\beta$。
\end{cor}

\begin{proof}
  令 $\beta = 2^\alpha$。
  现在我们要证明的是一个命题，因此由归纳法，可以假设 $\alpha$ 是 $\cd{A}$，于是 $\beta\jdeq \cd{A\to \bool}$。
  利用排中律，定义函数 $f:A\to (A\to \bool)$ 如下：
  \[f(a)(a') \defeq
  \begin{cases}
    \btrue &\quad a=a'\\
    \bfalse &\quad a\neq a'.
  \end{cases}
  \]
  如果 $f(a)=f(a')$，则 $f(a')(a) = f(a)(a) = \btrue$，从而 $a=a'$；因此 $f$ 是单射。
  所以，$\alpha \jdeq \cd{A} \le \cd{A\to \bool} \jdeq 2^\alpha$。

  另一方面，如果 $2^\alpha \le \alpha$，则会存在一个单射 $(A\to\bool)\to A$。
  由 \cref{thm:injsurj}，由于我们有 $(\lam{x} \bfalse):A\to \bool$ 且假设了排中律，则会存在一个满射 $A \to (A\to \bool)$，这与 Cantor 定理矛盾。
\end{proof}

\section{序数}
\label{sec:ordinals}

\index{ordinal|(}%

\begin{defn}\label{defn:accessibility}
  设 $A$ 是一个集合，且
  \[(\blank<\blank):A\to A\to \prop\]
  是 $A$ 上的一个二元关系。
  我们归纳地定义元素 $a:A$ 关于 $<$ 是\define{可达的}的含义：
  \indexdef{accessibility}%
  \indexsee{accessible}{accessibility}%
  \begin{itemize}
  \item 如果每个 $b<a$ 都是可达的，那么 $a$ 就是可达的。
  \end{itemize}
  我们记 $\acc(a)$ 表示 $a$ 是可达的。
\end{defn}

乍看之下，这样的归纳定义似乎永远无法启动，但显然，如果 $a$ 满足\emph{不存在} $b$ 使得 $b<a$ 这一性质，那么 $a$ 就是空虚地可达的。

注意，这是一个类型族的归纳定义，类似于 \cref{sec:generalizations} 中考虑的向量类型。
更准确地说，它有一个构造子，比如说 $\acc_<$，其类型为
\[ \acc_< : \prd{a:A} \Parens{\prd{b:A} (b<a) \to \acc(b)} \to \acc(a). \]
\index{induction principle!for accessibility}%
$\acc$ 的归纳原理是指，对于任意 $P:\prd{a:A} \acc(a) \to \type$，如果我们有
\[f:\prd{a:A}{h:\prd{b:A} (b<a) \to \acc(b)}
\Parens{\prd{b:A}{l:b<a} P(b,h(b,l))} \to
P(a,\acc_<(a,h)),
\]
那么就有一个由归纳定义的 $g:\prd{a:A}{c:\acc(a)} P(a,c)$，满足
\[g(a,\acc_<(a,h)) \jdeq f(a,\,h,\,\lam{b}{l} g(b,h(b,l))).\]
这一表述颇为繁琐，但我们通常只在 $P:A\to\type$ 仅依赖于 $A$ 这一更简单的情况下应用它。
在这种情况下，$f$ 的第二和第三个参数可以合并，因此我们需要证明的是：
\[f:\prd{a:A} \Parens{\prd{b:A} (b<a) \to \acc(b) \times P(b)}
\to P(a).
\]
也就是说，假设每个 $b<a$ 都是可达的且 $g(b):P(b)$ 已定义，由此定义 $g(a):P(a)$。

省略 $P$ 的第二个参数是合理的，这由下面的引理保证，而该引理的证明也是我们使用归纳原理更一般形式的唯一场合。

\begin{lem}
  可达性\index{accessibility}是一个命题。
\end{lem}

\begin{proof}
  我们必须证明，对于任意 $a:A$ 和 $s_1,s_2:\acc(a)$，有 $s_1=s_2$。
  我们通过对 $s_1$ 的归纳来证明这一点，令
  \[P_1(a,s_1) \defeq \prd{s_2:\acc(a)} (s_1=s_2). \]
  因此，我们必须证明，对于任意 $a:A$、${h_1:\prd{b:A} (b<a) \to \acc(b)}$ 以及
  \[ k_1:{\prd{b:A}{l:b<a}{t:\acc(b)} h_1(b,l) = t},\]
  对任意 $s_2:\acc(a)$ 都有 $\acc_<(a,h) = s_2$。
  我们将此陈述视为 $\prd{a:A}{s_2:\acc(a)} P_2(a,s_2)$，其中
  \[P_2(a,s_2) \defeq
  \prd{h_1 : \cdots } %{h_1:\prd{b:A} (b<a) \to \acc(b)}
  {k_1 : \cdots} % \Parens{\prd{b:A}{l:b<a}{t:\acc(b)} h_1(b,l) = t} \to
  (\acc_<(a,h_1) = s_2);
  \]
  因此可以通过对 $s_2$ 的归纳来证明。
  于是，假设 $h_2 : \prd{b:A} (b<a) \to \acc(b)$，以及 $k_2$ 具有一个极其复杂但在此无关紧要的类型，
  % \begin{narrowmultline*}
  %   k_2:\prd{b:A}{l:b<a}
  %   \prd{h_1:\prd{b':A} (b'<b) \to \acc(b')}
  %   \narrowbreak
  %   \Parens{\prd{b':A}{l':b'<b}{t':\acc(b')} h_1(b',l') = t'} \to
  %   (\acc_<(b,h_1) = h_2(b,l)).
  % \end{narrowmultline*}
  并且必须证明，对于任何具有上述类型的 $h_1$ 和 $k_1$，有 $\acc_<(a,h_1) = \acc_<(a,h_2)$。
  根据函数外延性，只需证明对所有 $b:A$ 和 $l:b<a$，有 $h_1(b,l) = h_2(b,l)$。
  而这由 $k_1$ 即得。
\end{proof}

\begin{defn}
  集合 $A$ 上的二元关系 $<$ 称为\define{良基的}
  \indexdef{relation!well-founded}%
  \indexdef{well-founded!relation}%
  如果 $A$ 的每个元素都是可达的。
\end{defn}

良基性的要旨在于：对于 $P:A\to \type$，我们可以利用 $\acc$ 的归纳原理得出 $\prd{a:A} \acc(a) \to P(a)$，然后应用良基性得出 $\prd{a:A} P(a)$。
换言之，如果从 $\fall{b:A} (b<a) \to P(b)$ 可以推出 $P(a)$，那么就有 $\fall{a:A} P(a)$。
这种方法称为\define{良基归纳}\indexdef{well-founded!induction}。

\begin{lem}
  良基性是一个命题。
\end{lem}

\begin{proof}
  $<$ 的良基性是类型 $\prd{a:A} \acc(a)$，由于每个 $\acc(a)$ 都是命题，所以它也是命题。
\end{proof}

\begin{eg}\label{thm:nat-wf}
  或许最常见的良基关系是 $\nat$ 上通常的严格序关系。
  为了证明它是良基的，需要对每个 $n:\nat$ 证明 $n$ 是可达的。
  \index{strong!induction}%
  这正是用 $\nat$ 上的普通归纳法证明"强归纳法"的常见证法。

  具体而言，我们对 $n:\nat$ 作归纳，证明对所有 $k\le n$，$k$ 都是可达的。
  基本情形是 $0$ 可达，这空虚成立，因为没有任何东西严格小于 $0$。
  在归纳步骤中，假设对所有 $k\le n$（即对所有 $k < n+1$），$k$ 都是可达的；于是根据定义，$n+1$ 也是可达的。

  $\nat$ 上的另一种良基关系可以通过如下方式定义：对所有 $n:\nat$，令 $n < \suc(n)$。
  这个关系的良基性几乎就是 $\nat$ 的普通归纳原理。
\end{eg}

\begin{eg}\label{thm:wtype-wf}
  设 $A:\set$ 且 $B : A \to \set$ 为一族集合。
  回顾 \cref{sec:w-types}，$W$ 类型 $\wtype{a:A} B(a)$ 由如下唯一的构造子归纳地生成：
  \begin{itemize}
  \item $\supp : \prd{a:A} (B(a) \to \wtype{x:A} B(x)) \to \wtype{x:A} B(x)$
  \end{itemize}
  我们通过对其第二个参数递归来定义 $\wtype{x:A} B(x)$ 上的关系 $<$：
  \begin{itemize}
  \item 对任意 $a:A$ 和 $f:B(a) \to \wtype{x:A} B(x)$，定义 $w<\supp(a,f)$ 为：仅仅存在 $b:B(a)$ 使得 $w = f(b)$。
  \end{itemize}
  现在我们用 $\wtype{x:A}B(x)$ 的通常归纳原理证明每个 $w:\wtype{x:A} B(x)$ 对此关系都是可达的。
  即假设给定 $a:A$ 和 $f:B(a) \to \wtype{x:A} B(x)$，以及一个提升 $f' : \prd{b:B(a)} \acc(f(b))$。
  于是根据 $<$ 的定义，对所有 $w<\supp(a,f)$ 都有 $\acc(w)$；因此 $\supp(a,f)$ 是可达的。
\end{eg}

良基性允许我们递归地定义函数、通过归纳证明命题，如下面的引理所示。
回顾 \cref{subsec:prop-subsets}，$\power B$ 表示\emph{幂集}\index{power set} $\power B \defeq (B\to\prop)$。

\begin{lem}\label{thm:wfrec}
  设 $B$ 为集合，且我们有一个函数
  \[ g : \power B \to B \]
  那么，如果 $<$ 是 $A$ 上的良基关系，则存在一个函数 $f:A\to B$，使得对所有 $a:A$ 都有
  \begin{equation*}
    f(a) = g\Big(\setof{ f(a') | a'<a }\Big).
  \end{equation*}
\end{lem}


\noindent
（这里使用了 \cref{sec:image} 中子集像的记号。）


\begin{proof}
  我们首先对每个 $a:A$ 和 $s:\acc(a)$ 定义一个元素 $\bar f(a,s):B$。
  由归纳原理，只需假设 $s$ 是一个函数，为每个 $a'<a$ 指派一个见证 $s(a'):\acc(a')$，而且对每个这样的 $a'$ 已有元素 $\bar f(a',s(a')):B$。
  在此情形下，我们定义
  \begin{equation*}
    \bar f(a,s) \defeq g\Big(\setof{ \bar f(a',s(a')) | a'<a }\Big).
  \end{equation*}

  现在，由于 $<$ 是良基的，我们有一个函数 $w:\prd{a:A} \acc(a)$。
  因此，我们可以定义 $f(a)\defeq \bar f (a,w(a))$。
\end{proof}

在经典\index{mathematics!classical}逻辑中，良基性有一种更为人熟知的等价表述。
在下文中，我们称一个子集 $B: \power A$ 是\define{非空的}
\indexdef{nonempty subset}
如果它与空子集 $(\lam{x}\bot) : \power X$ 不相等。
我们留待读者验证：在假设排中律的前提下，这等价于仅被居住，即条件 $\exis{x:A} x\in B$。

\begin{lem}\label{thm:wfmin}
  \index{excluded middle}%
  假设排中律，则 $<$ 是良基的当且仅当每个非空子集 $B: \power A$ 仅仅有一个极小元。
\end{lem}

\begin{proof}
  先假设 $<$ 是良基的，并设 $B\subseteq A$ 是一个没有极小元的子集。
  也就是说，对任意满足 $a\in B$ 的 $a:A$，都仅仅存在某个 $b:A$ 使得 $b<a$ 且 $b\in B$。

  我们断言：对任意 $a:A$ 和 $s:\acc(a)$，都有 $a\notin B$。对此作归纳证明。
  可假设 $s$ 是一个函数，对每个 $a'<a$ 指派一个证明 $s(a'):\acc(a')$，而且对每个这样的 $a'$，已有 $a'\notin B$。
  若 $a\in B$，则由假设，将仅仅存在某个 $b<a$ 使得 $b\in B$，这与归纳假设矛盾。
  因此，$a\notin B$；这就完成了归纳。
  由于 $<$ 是良基的，对所有 $a:A$ 都有 $a\notin B$，即 $B$ 是空集。

  现在反过来假设每个非空子集都仅仅有一个极小元。
  令 $B = \setof{ a:A | \neg \acc(a) }$。
  若 $B$ 非空，则它仅仅有一个极小元。
  于是，仅仅存在某个满足 $a\in B$ 的 $a:A$，使得对所有 $b<a$，都有 $\acc(b)$。
  但根据定义（并通过对截断作归纳），$a$ 仅仅可达，从而可达，这与 $a\in B$ 矛盾。
  因此 $B$ 是空集，故 $<$ 是良基的。
\end{proof}

\begin{defn}
  集合 $A$ 上的良基关系 $<$ 被称为\define{外延的}
  \indexdef{relation!extensional}%
  \indexdef{extensional!relation}%
  若对任意 $a,b:A$，都有
  \[ \Parens{\fall{c:A} (c<a) \Leftrightarrow (c<b)} \to (a=b). \]
\end{defn}

注意，由于 $A$ 是集合，外延性是命题。
这里的"外延性"与函数外延性无关，也与相等类型的外延性无关。
\index{axiom!of extensionality}%
更确切地说，它是经典集合论中外延公理的"局部"对应物。

\begin{thm}
  外延良基关系所构成的类型是一个集合。
\end{thm}

\begin{proof}
  根据泛等公理，只需证明：如果 $(A,<)$ 是外延且良基的，并且 $f:(A,<) \cong (A,<)$，那么 $f=\idfunc[A]$。
  \index{automorphism!of extensional well-founded relations}%
  我们通过对 $<$ 作归纳来证明对所有 $a:A$ 都有 $f(a)=a$。
  归纳假设是：对于所有 $a'<a$，有 $f(a')=a'$。

  现在，由于 $A$ 是外延的，要证 $f(a)=a$，只需证明
  \[\fall{c:A}(c<f(a)) \Leftrightarrow (c<a).\]
  然而，因为 $f$ 是一个自同构，我们有 $(c<a) \Leftrightarrow (f(c)<f(a))$。
  但由归纳假设，$c<a$ 意味着 $f(c)=c$，所以 $(c<a) \to (c<f(a))$。
  另一方面，如果 $c<f(a)$，那么 $f^{-1}(c)<a$，从而再次利用归纳假设得到 $c = f(f^{-1}(c)) = f^{-1}(c)$；于是有 $c<a$。
  因此，对任意 $c:A$，我们有 $(c<a) \Leftrightarrow (c<f(a))$，故 $f(a)=a$。
\end{proof}

\begin{defn}\label{def:simulation}
  如果 $(A,<)$ 和 $(B,<)$ 都是外延且良基的，则称满足以下条件的函数 $f:A\to B$ 为\define{模拟}
  \indexdef{simulation}%
  \indexsee{function!simulation}{simulation}%
  ：
  \begin{enumerate}
  \item 若 $a<a'$，则 $f(a)<f(a')$；并且\label{item:sim1}
  \item 对所有 $a:A$ 和 $b:B$，若 $b<f(a)$，则仅仅存在 $a'<a$ 满足 $f(a')=b$。\label{item:sim2}
  \end{enumerate}
\end{defn}

\begin{lem}
  任意模拟都是单的。
\end{lem}

\begin{proof}
  我们通过双重良基归纳法证明：对任意 $a,b:A$，若 $f(a)=f(b)$，则 $a=b$。
  关于 $a:A$ 的归纳假设是：对任意 $a'<a$ 和任意 $b:B$，若 $f(a')=f(b)$，则 $a=b$。
  关于 $b:A$ 的内层归纳假设是：对任意 $b'<b$，若 $f(a)=f(b')$，则 $a=b'$。

  假设 $f(a)=f(b)$；我们需要证明 $a=b$。
  由外延性，只需证明对任意 $c:A$ 我们有 $(c<a)\Leftrightarrow (c<b)$。
  如果 $c<a$，那么由 \cref{def:simulation}\ref{item:sim1} 可得 $f(c)<f(a)$。
  故 $f(c)<f(b)$，于是由 \cref{def:simulation}\ref{item:sim2} 可知，仅仅存在 $c':A$ 使得 $c'<b$ 且 $f(c)=f(c')$。
  由关于 $a$ 的归纳假设，我们有 $c=c'$，因此 $c<b$。
  对偶的论证同理可证。
\end{proof}

特别地，这表明在 \cref{def:simulation}\ref{item:sim2} 中，"仅仅"一词可以省略而不改变意思。

\begin{cor}
  如果 $f:A\to B$ 是模拟，那么对所有 $a:A$ 和 $b:B$，若 $b<f(a)$，则\emph{纯粹存在} $a'<a$ 使得 $f(a')=b$。
\end{cor}

\begin{proof}
  由于 $f$ 是单的，$\sm{a:A} (f(a)=b)$ 是命题。
\end{proof}

我们称子集 $C :\power B$ 为\define{初始段}
\indexdef{initial!segment}%
\indexsee{segment, initial}{initial segment}%
，如果 $c\in C$ 且 $b<c$ 蕴含 $b\in C$。
模拟的像必为初始段，而任意初始段的包含映射都是模拟。
因此，由泛等性，每个模拟 $A\to B$ 都\emph{相等于} $B$ 的某个初始段的包含映射。

\begin{thm}
  对一个集合 $A$，令 $P(A)$ 为 $A$ 上外延良基关系的类型。
  若 $\mathord{<_A} : P(A)$ 与 $\mathord{<_B} : P(B)$ 且 $f:A\to B$，令 $H_{\mathord{<_A}\mathord{<_B}}(f)$ 为命题"$f$ 是模拟"。
  则 $(P,H)$ 是 \cref{sec:sip} 意义下 \uset 上的标准结构概念。
\end{thm}

\begin{proof}
  恒等映射是模拟、模拟的复合仍是模拟，这些留给读者验证。
  因此，模拟构成了一个结构概念。
  为验证标准性，我们须证：若 $<$ 和 $\prec$ 是 $A$ 上的两个外延良基关系，且 $\idfunc[A]$ 在两个方向都是模拟，则 $<$ 与 $\prec$ 相等。
  由于外延性和良基性都是命题，为此只需 $\fall{a,b:A} (a<b) \Leftrightarrow (a\prec b)$ 成立。
  而这由 $\idfunc[A]$ 的 \cref{def:simulation}\ref{item:sim1} 即可推出。
\end{proof}

\begin{cor}\label{thm:wfcat}
  存在一个范畴，其对象为配备外延良基关系的集合，态射为模拟。
\end{cor}

实际上，这个范畴是一个偏序集。

\begin{lem}
  对于外延良基的 $(A,<)$ 和 $(B,<)$，至多存在一个模拟 $f:A\to B$。
\end{lem}

\begin{proof}
  假设 $f,g:A\to B$ 都是模拟。
  由于"是模拟"是一个命题，只需证明 $\fall{a:A}(f(a)=g(a))$。
  对 $<$ 作归纳，可假设对所有 $a'<a$ 均有 $f(a')=g(a')$。
  而由 $B$ 的外延性，要证 $f(a)=g(a)$，只需 $\fall{b:B}(b<f(a)) \Leftrightarrow (b<g(a))$ 成立。

  由于 $f$ 是模拟，若 $b<f(a)$，则存在 $a'<a$ 使得 $f(a')=b$。
  由归纳假设亦有 $g(a')=b$，从而 $b<g(a)$。
  对偶论证同理可证。
\end{proof}

因此，若 $A$ 和 $B$ 配备了外延良基关系，可记 $A\le B$ 表示存在模拟 $f:A\to B$。
\cref{thm:wfcat} 表明：若 $A\le B$ 且 $B\le A$，则 $A=B$。

\begin{defn}
  一个\define{序数}
  \indexdef{ordinal}%
  \indexsee{number!ordinal}{ordinal}%
  是一个集合 $A$ 连同其上的一个外延良基关系，且该关系是\emph{传递的}，即满足 $\fall{a,b,c:A}(a<b)\to (b<c) \to (a<c)$。
\end{defn}

\begin{eg}
  当然，$\nat$ 上通常的严格序是传递的。
  不难看出它也是外延的；因此它是一个序数。
  依照惯例，我们记这个序数为 $\omega$。
\end{eg}

\symlabel{ord}
令 $\ord$ 表示序数的类型。
根据之前的结果，$\ord$ 是一个集合并具有自然的偏序。
我们现在证明 $\ord$ 也容许一个良基关系。

\symlabel{initial-segment}
如果 $A$ 是一个序数且 $a:A$，令 $\ordsl A a \defeq \setof{ b:A | b<a}$ 表示\emph{初始段}。
\index{initial!segment}%
注意，若 $\ordsl A a = \ordsl A b$ 作为序数相等，则该同构必须保持它们到 $A$ 的包含映射（因为模拟构成偏序集），从而它们作为 $A$ 的子集也是相等的。
因此，由 $A$ 的外延性知 $a=b$。
于是，函数 $a\mapsto \ordsl A a$ 是一个单射 $A\to \ord$。

\begin{defn}
  对于序数 $A$ 和 $B$，一个模拟 $f:A\to B$ 被称为\define{有界的}
  \indexdef{simulation!bounded}%
  \indexdef{bounded!simulation}%
  如果存在 $b:B$ 使得 $A = \ordsl B b$。
\end{defn}

上述讨论表明，这样的 $b$ 若存在则唯一，因此有界性是一个命题。

我们记 $A<B$ 表示存在一个从 $A$ 到 $B$ 的有界模拟。
由于模拟是唯一的，$A<B$ 也是一个命题。

\begin{thm}\label{thm:ordord}
  $(\ord,<)$ 是一个序数。
\end{thm}

\noindent
更准确地说，该定理断言：某个宇宙\index{universe level}中的序数类型 $\ord_{\UU_i}$ 本身是更高一层宇宙中的序数，即 $(\ord_{\UU_i},<):\ord_{\UU_{i+1}}$。

\begin{proof}
  设 $A$ 是一个序数；首先证明对所有 $a:A$，$\ordsl A a$ 在 $\ord$ 中是可达的。
  对 $A$ 作良基归纳，假设对所有 $b<a$，$\ordsl A b$ 是可达的。
  根据可达性的定义，须证对所有 $B<\ordsl A a$，$B$ 在 $\ord$ 中是可达的。
  然而，若 $B<\ordsl A a$，则存在某个 $b<a$ 使得 $B = \ordsl{(\ordsl A a)}{b} = \ordsl A b$，而由归纳假设知后者是可达的。
  因此，对所有 $a:A$，$\ordsl A a$ 是可达的。

  现在证明 $A$ 在 $\ord$ 中是可达的。根据定义，须证对所有 $B<A$，$B$ 是可达的。
  但如前所述，$B<A$ 意味着存在某个 $a:A$ 使得 $B=\ordsl A a$，而由刚才的证明知这是可达的。
  因此，$\ord$ 是良基的。

  对于外延性，假设 $A$ 和 $B$ 是满足下式的序数：
  \narrowequation{\prd{C:\ord} (C<A) \Leftrightarrow (C<B).}
  那么对每个 $a:A$，由于 $\ordsl A a<A$，我们有 $\ordsl A a<B$，从而存在 $b:B$ 使得 $\ordsl A a = \ordsl B b$。
  定义 $f:A\to B$ 将每个 $a$ 映到相应的 $b$；可以直接验证 $f$ 是一个同构。
  于是 $A\cong B$，由泛等性知 $A=B$。

  最后，容易看出 $<$ 是传递的。
\end{proof}

将 $\ord$ 视为序数通常非常方便，但也有其隐患。
例如，考虑以下引理，我们需留意其中宇宙的使用方式。

\begin{lem}\label{thm:ordsucc}
  设 \bbU 为一个宇宙。
  对于任意 $A:\ord_\bbU$，存在 $B:\ord_\bbU$ 使得 $A<B$。
\end{lem}

\begin{proof}
  令 $B=A+\unit$，并规定 $\unit$ 中的元素 $\ttt$ 大于 $A$ 的所有元素。
  则 $B$ 是一个序数，且容易看出 $A\cong \ordsl B \ttt$。
\end{proof}

在 \cref{thm:ordsucc} 证明中所构造的序数 $B$ 称为 $A$ 的\define{后继}\indexdef{successor!of an ordinal}。

这条引理揭示了使用"类型歧义"\index{typical ambiguity}风格——即用 \UU 表示一个任意且未指定的宇宙——可能带来的隐患。
考虑如下另一种证明。

\begin{proof}[\cref{thm:ordsucc} 的另一种拟议证明]
  注意到 $C<A$ 当且仅当存在某个 $a:A$ 使得 $C=\ordsl A a$。
  这给出了一个同构 $A \cong \ordsl \ord A$，从而 $A<\ord$。
  于是我们可以取 $B\defeq\ord$。
\end{proof}

如果我们用类型歧义的风格陈述 \cref{thm:ordsucc}，则第二种证明是有效的。
但这样得到的引理用处较小，因为第二种证明会强制引理陈述中第二个"\ord"所指的宇宙层级高于第一个"\ord"。
而第一种证明则允许两个宇宙是同一个。

类似的讨论也适用于下一条引理：该引理也可通过观察"对任意 $A:\ord$ 皆有 $A\le \ord$"来证明，但那样得到的结论用处较小。

\begin{lem}\label{thm:ordunion}
  设 \bbU 为一个宇宙。
  对于任意 $X:\type$ 与 $F:X\to \ord_\bbU$，存在 $B:\ord_\bbU$ 使得对所有 $x:X$ 均有 $Fx\le B$。
\end{lem}

\begin{proof}
  在 $\sm{x:X} Fx$ 上定义等价关系 $\eqr$ 如下，并令 $B$ 为其集合商（参见 \cref{rmk:quotient-of-non-set}）：
  \[ (x,y) \eqr (x',y')
  \;\defeq\;
  \Big(\ordsl{(Fx)}{y} \cong \ordsl{(Fx')}{y'}\Big).
  \]
  定义 $(x,y)<(x',y')$ 当且仅当 $\ordsl{(Fx)}{y} < \ordsl{(Fx')}{y'}$。
  该关系显然下降到商上，并可验证使 $B$ 成为一个序数。
  此外，对每个 $x:X$，诱导的映射 $Fx\to B$ 是一个模拟。
\end{proof}

\section{经典良序}
\label{sec:wellorderings}

\index{denial|(}%
现在我们证明，我们的序数与更为人熟知的经典\index{mathematics!classical}良序等价。

\begin{lem}
  \index{excluded middle}%
  假设排中律，每个序数都是三分的：
  \index{trichotomy of ordinals}%
  \index{ordinal!trichotomy of}%
  \[ \fall{a,b:A} (a<b) \vee (a=b) \vee (b<a). \]
\end{lem}

\begin{proof}
  对 $a$ 归纳，可假设对任意 $a'<a$ 及任意 $b':A$，有 $(a'<b') \vee (a'=b') \vee (b'<a')$。
  现在对 $b$ 归纳，可假设对任意 $b'<b$，有 $(a<b') \vee (a=b') \vee (b'<a)$。

  根据排中律，要么仅存在 $b'<b$ 使得 $a<b'$，要么仅存在 $b'<b$ 使得 $a=b'$，要么对任意 $b'<b$ 皆有 $b'<a$。
  在第一种情形，由传递性仅知 $a<b$，而 $a<b$ 是命题，故确实有 $a<b$。
  类似地，在第二种情形，由迁移可得 $a<b$。
  因此，我们假设 $\fall{b':A}(b'<b)\to (b'<a)$。

  现在类似地，要么仅存在 $a'<a$ 使得 $b<a'$，要么仅存在 $a'<a$ 使得 $a'=b$，要么对任意 $a'<a$ 皆有 $a'<b$。
  在第一种和第二种情形，$b<a$，因此我们可假设 $\fall{a':A}(a'<a)\to (a'<b)$。
  然而，由外延性，这两个假设现在蕴涵 $a=b$。
\end{proof}

\begin{lem}
  良基关系不包含循环，即
  \[ \fall{n:\mathbb{N}}{a:\mathbb{N}_n\to A} \neg\Big((a_0<a_1) \wedge \dots \wedge (a_{n-1}<a_n)\wedge (a_n<a_0)\Big). \]
\end{lem}

\begin{proof}
  我们对 $a:A$ 归纳证明：不存在包含 $a$ 的循环。
  于是，由归纳可假设对所有 $a'<a$，不存在包含 $a'$ 的循环。
  但在任何包含 $a$ 的循环中，都存在某个小于 $a$ 且属于该循环的元素。
\end{proof}

\indexdef{relation!irreflexive}%
\index{irreflexivity!of well-founded relation}%
特别地，良基关系必须是\define{不自反的}，即对任意 $a$ 有 $\neg(a<a)$。

\begin{thm}\label{thm:wellorder}
  假设排中律，则 $(A,<)$ 是一个序数当且仅当 $A$ 的每个非空子集 $B\subseteq A$ 都有最小元。
\end{thm}

\begin{proof}
  若 $A$ 是序数，则由 \cref{thm:wfmin}，其每个非空子集都仅有一个极小元。
  但三分性保证了任何极小元都是最小元。
  此外，最小元若存在则唯一，因此仅存在一个最小元与实际拥有一个效果相同。

  反之，若每个非空子集都有最小元，则由 \cref{thm:wfmin} 知 $A$ 是良基的。
  我们也有三分性，因为对任意 $a,b$，子集
  $ \setof{a,b} \defeq \setof{x:A | x=a \lor x=b} $
  仅有一个最小元，而该最小元必定是 $a$ 或 $b$。
  这蕴含传递性：若 $a<b$ 且 $b<c$，则 $a=c$ 或 $c<a$ 都会产生循环。
  类似地，它也蕴含外延性：若 $\fall{c:A}(c<a)\Leftrightarrow (c<b)$，则 $a<b$ 将导致（取 $c$ 为 $a$ 时）$a<a$，这构成循环；$b<a$ 的情形同理；故必有 $a=b$。
\end{proof}

在经典\index{mathematics!classical}数学中，\cref{thm:wellorder} 这一刻画被直接当作\emph{良序}的定义，而\emph{序数}则构成良序同构类的一组标准代表元。
在我们的语境中，结构同一性原理意味着无需寻找这样的代表元：任何一个良序都不逊于任何其他良序。

现在我们转而考虑选择公理的推论。
对任意集合 $X$，令 $\powerp X$ 表示由 $X$ 的\emph{仅知有实例的子集}构成的类型：
\symlabel{inhabited-powerset}
\[ \powerp X \defeq \setof{ Y : \power X | \exis{x:X} x\in Y}. \]
若假设排中律，这便等价于由 $X$ 的\emph{非空}\index{nonempty subset}子集构成的类型，且有 $\power X \eqvsym (\powerp X) + \unit$。

\begin{thm}\label{thm:wop}
  \index{axiom!of choice}%
  \index{excluded middle}%
  在假设排中律的前提下，以下两条陈述等价：
  \begin{enumerate}
  \item 对每个集合 $X$，仅存在一个函数 $f: \powerp X \to X$，使得对所有 $Y:\powerp X$ 都有 $f(Y)\in Y$。\label{item:wop1}
  \item 每个集合仅具有一个序数结构。\label{item:wop2}
  \end{enumerate}
\end{thm}

\noindent
当然，~\ref{item:wop1} 是选择公理的经典\index{mathematics!classical}标准形式；参见 \cref{ex:choice-function}。

\begin{proof}
  一个方向的证明很简单：假设~\ref{item:wop2}。
  由于我们要证明的是命题~\ref{item:wop1}，可以假定 $A$ 是一个序数。
  于是就可以把 $f(B)$ 定义为 $B$ 中的最小元。

  现在假设~\ref{item:wop1}。
  同前，由于~\ref{item:wop2} 是一个命题，可以假定已经给定了这样的 $f$。
  我们以显然的方式将 $f$ 延拓为函数
  \[ \bar f:\power X \eqvsym (\powerp X) + \unit \longrightarrow X+\unit
  \]

  现在，对任意序数 $A$，可通过良基递归定义 $g_A : A \to X + \unit$：
  \[ g_A(a) \defeq
    \bar f\Big(X \setminus \setof{ g_A(b) | \strut (b<a) \wedge (g_A(b) \in X) }\Big)
  \]
  （显然地将 $X$ 视作 $X + \unit$ 的子集）。

  令 $A'\defeq \setof{a:A | g_A(a) \in X}$ 为 $X \subseteq X + \unit$ 的原像；则我们断言限制映射 $g_A' : A' \to X$ 是单射。
  因为若 $a, a' : A$ 且 $a \neq a'$，则由三分性且不失一般性，可假设 $a' < a$。
  于是 $g_A(a') \in \setof{ g_A(b) | b<a }$，再由对所有 $Y$ 都有 $f(Y) \in Y$ 便知 $g_A(a) \neq g_A(a')$。

  此外，$A'$ 是 $A$ 的一个初始段。
  因为 $g_A(a)$ 属于 \unit 当且仅当 $\setof{g_A(b)|b<a} = X$，且若此条件成立，则它对任意 $a' > a$ 亦成立。
  因此，$A'$ 本身也是一个序数。

  最后，由于 \ord 是一个序数，我们可以取 $A \defeq \ord$。
  令 $X'$ 为 $g_\ord' : \ord' \to X$ 的像；则 $g_\ord'$ 的逆给出一个单射 $H : X' \to \ord$。
  由 \cref{thm:ordunion}，存在一个序数 $C$，使得对所有 $x : X'$ 都有 $Hx \le C$。
  再由 \cref{thm:ordsucc}，存在另一个序数 $D$ 使得 $C < D$，从而对所有 $x : X'$ 都有 $Hx < D$。
  现在我们有
  \begin{align*}
    g_{\ord}(D) &= \bar f\Big( X \setminus \setof{ g_\ord(B) | \rule{0pt}{1em} B<D \wedge (g_\ord(B) \in X)} \Big)\\
    &=\bar f\Big( X \setminus \setof{ g_\ord(B) | \rule{0pt}{1em} g_\ord(B) \in X} \Big)
  \end{align*}
  因为若 $B : \ord$ 且 $g_\ord(B) \in X$，则存在某个 $x : X'$ 使得 $B = Hx$，从而 $B < D$。
  现在，如果
  \[\setof{ g_\ord(B) | \rule{0pt}{1em} g_\ord(B) \in X}\]
  不等于整个 $X$，则 $g_\ord(D)$ 将属于 $X$ 却不在此子集中，而由于 $D$ 本身也是 $B$ 的一个可能取值，这将导致矛盾。
  因此这个子集必须是整个 $X$，从而 $g_\ord'$ 不仅是单射，还是满射。
  于是，我们就可以将 $\ord'$ 上的序数结构迁移到 $X$ 上。
\end{proof}

\begin{rmk}
  如果我们之前对 \cref{thm:ordsucc} 或 \cref{thm:ordunion} 给出了错误的证明，那么由此得到的 \cref{thm:wop} 的证明将是无效的：因为那样将无法一致地指定宇宙层级\index{universe level}。
  事实上，我们需要命题大小重整（它可由 \LEM{} 推出）来确保 $X'$ 在等价意义下与 $X$ 处于同一宇宙。
\end{rmk}

\begin{cor}
  假设选择公理，函数 $\ord\to\set$（它忘记序结构）是一个满射。
\end{cor}

注意，\ord 是一个集合，而 \set 是一个 $1$-类型。
一般而言，没有理由指望 $1$-类型能接受来自某个集合的满射函数。
即使选择公理似乎也不蕴含\emph{每个} $1$-类型都能如此（尽管参见 \cref{ex:acnm-surjset}），但它可以直接推出这一点对于由 \set 构造出来的 $1$-类型是成立的，例如 \cref{sec:sip} 中结构范畴的对象类型。
下面的推论也适用于这类范畴。

\begin{cor}
  \index{weak equivalence!of precategories}%
  假设 \choice{}，\uset 接受一个来自严格范畴的弱等价函子。
\end{cor}

\begin{proof}
  令 $X_0\defeq \ord$，并且对于 $A,B:X_0$，令 $\hom_X(A,B) \defeq (A\to B)$。
  由于 \ord 是一个集合，$X$ 是一个严格范畴，且上述满射 $X_0 \to \set$ 可扩张为弱等价函子 $X\to \uset$。
\end{proof}

现在回顾 \cref{sec:cardinals}，我们还有一个进一步的满射 $\cd{\blank}:\set\to\card$，从而得到一个复合的满射 $\ord\to\card$，它把每个序数映到它的基数。

\begin{thm}
  假设 \choice{}，满射 $\ord\to\card$ 具有一个截面。
\end{thm}

\begin{proof}
  对此有一个简单但错误的证明：因为 \ord 和 \card 都是集合，\choice{} 蕴含它们之间的任何满射\emph{仅}存在一个截面。
  然而，我们实际上拥有一个\emph{具体的}典范截面：因为 \ord 是一个序数，其每个非空子集都有一个唯一确定的最小元。
  因此，我们可以将每个基数映射到其对应纤维中的那个最小元。
\end{proof}

集合论中的传统做法是将基数等同于它们在 \ord 中的像，即具有该基数的那个最小序数。

由此可知，\card 也典范地具有序数结构：事实上，是一个与 \ord 同构的结构。
具体来说，我们通过良基递归定义一个函数 $\aleph:\ord\to\ord$，使得 $\aleph(A)$ 是满足如下条件的最小序数：其基数大于所有 $a:A$ 对应的 $\aleph({\ordsl A a})$。
那么（在假设 \choice{} 下）$\aleph$ 的像恰好就是 \card 的像。

\index{denial|)}%

\index{ordinal|)}%

\section{累积层次}
\label{sec:cumulative-hierarchy}

\index{bargaining|(}%
对于给定的宇宙 $\UU$，我们可以将其中所有集合的累积层次 $V$ 定义为一个高阶归纳类型。这种定义方式使得 $V$ 本身也是一个集合（位于更大的宇宙 $\UU'$ 中），并配备一个满足集合论常见定律的二元“隶属”关系 $x\in y$。

\begin{defn}\label{defn:V}
  相对于类型宇宙 $\UU$ 的\define{累积层次}
  \indexdef{cumulative!hierarchy, set-theoretic}%
  \indexdef{hierarchy!cumulative, set-theoretic}%
  $V$ 是由以下构造子生成的高阶归纳类型。
  %
  \begin{enumerate}
  \item 对于每一个 $A : \UU$ 和 $f : A \to V$，有一个元素 $\vset(A, f) : V$。
  \item 对所有 $A, B : \UU$，以及 $f : A \to V$ 和 $g : B \to V$，若满足
    %
    \begin{narrowmultline} \label{eq:V-path}
      \big(\fall{a:A} \exis{b:B} \id[V]{f(a)}{g(b)}\big) \land \narrowbreak
      \big(\fall{b:B} \exis{a:A} \id[V]{f(a)}{g(b)}\big)
    \end{narrowmultline}
    %
    则存在一条路径 $\id[V]{\vset(A,f)}{\vset(B,g)}$。
  \item 0-截断构造子：对所有 $x,y:V$ 和 $p,q:x=y$，有 $p=q$。
  \end{enumerate}
\end{defn}

用集合论的语言来说，$\vset(A,f)$ 可以理解为经典集合论意义下 $A$ 在 $f$ 下的像所构成的集合，即 $\setof{ f(a) | a \in A }$。
但是，我们将避免使用这种记法，因为它会与我们用于子类型的记法冲突（不过可以参见下面的 \eqref{eq:class-notation} 和 \cref{def:TypeOfElements}）。

层次 $V$ 是从空映射 $\rec\emptyt(V) : \emptyt \to V$（由空类型的递归原理给出）起步的，它给出的空集是 $\emptyset = \vset(\emptyt,\rec\emptyt(V))$。
接着，单点集 $\{\emptyset\}$ 通过映射 $\unit \to V$（定义为 $\ttt \mapsto \emptyset$）进入 $V$，以此类推。
（通过添加额外的点构造子，该定义也可调整以包含任意由“原子”或“本元”构成的集合。）
类型 $V$ 与基础宇宙 $\UU$ 处于同一宇宙层级。

$V$ 的第二个构造子的形式与我们之前见过的任何构造子都不同：它不仅仅涉及 $V$ 中的路径（我们在 \cref{sec:hittruncations} 中曾说过这有点奇怪），还涉及对其求和的截断。
它当然不符合 \cref{sec:naturality} 中描述的一般模式，因此它的归纳原理应该是什么可能并不显然。
幸运的是，就像我们在 \cref{sec:hittruncations} 中对 0-截断的第一个定义那样，它可以利用辅助的高阶归纳类型重新表达。
我们将其细节留给读者探索（参见 \cref{ex:cumhierhit}）。

\index{induction principle!for cumulative hierarchy}%
最终，$V$ 的归纳原理（用模式匹配的语言写出来）是说：给定 $P:V\to \set$，为了构造 $h:\prd{x:V} P(x)$，只需给出以下内容。

\begin{enumerate}
\item 对于任意 $f:A\to V$，在假设已对所有 $a:A$ 给出 $h(f(a))$ 的前提下，构造 $h(\vset(A,f))$。
\item 验证若 $f : A \to V$ 和 $g : B \to V$ 满足~\eqref{eq:V-path}，则 $\dpath{P}{q}{h(\vset(A,f))}{h(\vset(B,g))}$，其中 $q$ 是由 $V$ 的第二个构造子及~\eqref{eq:V-path} 产生的路径；这里需归纳地假设 $h(f(a))$ 和 $h(g(b))$ 对所有 $a:A$ 和 $b:B$ 均有定义，并且满足以下条件：
\begin{eqnarray*}
    &       & \big(\fall{a:A} \exis{b:B} \exis{p:f(a)=g(b)} \dpath{P}{p}{h(f(a))}{h(g(b))}\big) \\
    & \land & \big(\fall{b:B} \exis{a:A} \exis{p:f(a)=g(b)} \dpath{P}{p}{h(f(a))}{h(g(b))}\big)
\end{eqnarray*}
\end{enumerate}

第二个条款检查正在定义的映射必须尊重在 \eqref{eq:V-path} 中引入的路径。
正如我们通常在使用模式匹配陈述高阶归纳原理时那样，这看起来像是同义反复，其实不然。
关键在于，“$h(f(a))$”本质上是一个我们无法窥探其内部的形式符号，而 $h(\vset(A,f))$ 必须依据它来定义。
因此，在第二个条款中，我们在适当的时候假设这些形式符号相等，然后验证由第一个条款的构造所产生的元素也相等。
当然，如果 $P$ 是一个命题族，那么第二个条款会自动满足。

注意到，通过归纳可知，对于每个 $v:V$，仅有 $A:\UU$ 和 $f:A\to V$ 使得 $v=\vset(A,f)$。
因此，尝试如下定义 $V$ 上的 \define{隶属关系}
\indexdef{membership, for cumulative hierarchy}%
$x\in v$ 是合理的：
%
% Note: "membership" rather than "elementhood", because "element" is taken.
%
\symlabel{V-membership}
\begin{equation*}
  (x \in \vset(A,f)) \defeq (\exis{a : A} x = f(a)).
\end{equation*}
%
要验证该定义是有效的，我们必须使用 $V$ 的递归原理。假设我们有一条通过~\eqref{eq:V-path} 构造的道路 $\vset(A, f) = \vset(B, g)$。若 $x \in \vset(A,f)$，则仅有 $a : A$ 使得 $x = f(a)$；但由~\eqref{eq:V-path}，仅有 $b : B$ 使得 $f(a) = g(b)$，因此 $x = g(b)$ 且 $x \in \vset(B,g)$。反之亦然。

$V$ 上的 \define{子集关系}
\indexdef{subset!relation on the cumulative hierarchy}%
$x\subseteq y$ 按通常方式定义为：
%
\begin{equation*}
  (x \subseteq y) \defeq \fall{z : V} z \in x \Rightarrow z \in y.
\end{equation*}

一个 \define{类}
\indexdef{class}%
可以视作 $V$ 上的一个纯命题谓词。我们说一个类 $C : V \to \prop$ 是一个 \define{$V$-集（$V$-set）}
\indexdef{set!in the cumulative hierarchy}%
，如果仅存在 $v:V$ 使得
%
\begin{equation*}
  \fall{x : V} C(x) \Leftrightarrow x \in v.
\end{equation*}
我们也可以使用类的常规记法，这与我们对于子类型的标准记法一致：
\begin{equation}
  \setof{ x | C(x) } \defeq \lam{x}C(x).\label{eq:class-notation}
\end{equation}
%
一个类 $C: V\to \prop$ 被称为 \define{$\UU$-小的（$\UU$-small）}
\indexdef{class!small}%
\indexdef{small!class}%
，如果它的所有值 $C(x)$ 都位于 $\UU$ 中，具体而言即 $C: V\to \prop_{\UU}$。
由于 $V$ 与构建它所基于的基础宇宙 $\UU$ 位于同一宇宙 $\UU'$ 中，这对于任意 $v,w:V$ 的相等类型 $v=_V w$ 同样成立。因此，为了在缺乏命题大小重整的情况下获得一个行为良好的理论，
\index{propositional!resizing}%
\index{resizing}%
拥有一个 $\UU$-小的相等关系的“重整”会很方便，我们可以如下通过归纳来定义它。

\begin{defn}\label{def:bisimulation}
  通过对 $V$ 进行双重归纳来定义 \define{双模拟}
  \indexdef{bisimulation}%
  关系
  %
  \begin{equation*}
    \mathord\bisim : V \times V \longrightarrow \prop_{\UU}
  \end{equation*}
  %
  ，其中对于 $\vset(A,f)$ 和 $\vset(B,g)$，我们令：
  \begin{narrowmultline*}
    \vset(A,f)  \bisim \vset(B,g) \defeq \narrowbreak
    \big(\fall{a:A}\exis{b:B} f(a)  \bisim g(b)\big) \land
    \big(\fall{b:B}\exis{a:A} f(a) \bisim g(b)\big).
  \end{narrowmultline*}
\end{defn}


%
为验证该定义是正确的，我们只需检查它与通过~\eqref{eq:V-path} 构造的道路 $\vset(A, f) = \vset(B, g)$ 相容，而这是显然的；此外还需验证 $\prop_{\UU}$ 是一个集合，而它确实是。注意，由构造可知 $u \bisim v$ 位于 $\propU$ 中。

\begin{lem}\label{lem:BisimEqualsId}
对于任意 $u,v:V$，我们有 $(u=_V v) = (u \bisim v)$。
\end{lem}

\begin{proof}
容易通过归纳证明 $\bisim$ 是自反的，因此通过传输我们有 $(u=_V v)\to (u \bisim v)$。
因此，只需再证 $(u \bisim v)\to (u=_V v)$。
通过对 $u$ 和 $v$ 进行归纳，我们可以假设它们分别是 $\vset(A,f)$ 和 $\vset(B,g)$。
（我们可以忽略 $V$ 的道路构造子，因为 $(u \bisim v)\to (u=_V v)$ 是一个命题。）
那么根据定义，$\vset(A,f)\bisim\vset(B,g)$ 蕴含 $(\fall{a:A}\exis{b:B}f(a)  \bisim g(b))$，反之亦然。
而归纳假设告诉我们 $(\fall{a:A}\exis{b:B}f(a) = g(b))$，反之亦然。
因此，由 $V$ 的道路构造子，我们有 $\vset(A,f) =_V \vset(B,g)$。
\end{proof}

人们也许认为可以省略 $V$ 的 0-截断构造子，并通过将 \cref{thm:h-set-refrel-in-paths-sets} 应用于双模拟来\emph{证明} $V$ 是 0-截断的。然而，在 \cref{lem:BisimEqualsId} 的证明中，我们使用了 $V$ 是 0-截断的这一事实，从而得出 $(u \bisim v)\to (u=_V v)$ 是一个命题，使得在归纳中只需假设 $u$ 和 $v$ 是 $\vset(A,f)$ 和 $\vset(B,g)$ 即可。

现在，我们可以利用重整后的相等关系来得到以下有用的原理。

\begin{lem}\label{lem:MonicSetPresent}
对于每个 $u:V$，都存在给定的 $A_u:\UU$ 和单射 $m_u: A_u \mono V$，使得 $u = \vset(A_u, m_u)$。
\end{lem}

\begin{proof}
  任取一个表示 $u = \vset(A,f)$，并将 $f:A\to V$ 分解为一个满射后接一个单射：
  %
  \begin{equation*}
    f = m_u\circ e_u : A \epi A_u \mono V.
  \end{equation*}
  %
  显然 $u = \vset(A_u, m_u)$，只要 $A_u$ 仍在 $\UU$ 中即可；而若 $e_u : A \epi A_u$ 的核在 $\UU$ 中，则此条件成立。但 $e_u : A \epi A_u$ 的核是 $V$ 上恒等映射沿 $f : A\to V$ 的拉回，而由我们刚刚所证，在等价意义下它是 $\UU$-小的。现在，从 $u:V$ 构造满足 $m_u :A_u \mono V$ 和 $u = \vset(A_u, m_u)$ 的对 $(A_u, m_u)$ 在 $V$ 之上等价意义下唯一，从而由泛等性知在相等意义下唯一。于是，由唯一选择原理~\eqref{cor:UC}，存在映射 $c : V\to\sm{A:\UU}(A\to V)$ 使得 $c(u) = (A_u, m_u)$，其中 $m_u :A_u \mono V$ 且 $u = \vset(c(u))$，如所断言。
\end{proof}

\begin{defn}\label{def:TypeOfElements}
  对于 $u:V$，刚才构造的满足 $u = \vset(A_u, m_u)$ 的单的表示 $m_u: A_u \mono V$ 可以称为 $u$ 的 \define{成员类型}
  \indexdef{type!of members}%
  ，并记作 $m_u : [u] \mono V$，甚至简记为 $[u] \mono V$。我们可以将 $[u]$ 视作“由 $u$ 的成员组成的 $V$ 的子类”。
\end{defn}

\begin{thm}\label{thm:VisCST}
  \index{axiom!of set theory, for the cumulative hierarchy}%
  对于 $(V, {\in})$，以下性质成立：
  %
  \begin{enumerate}
  \item \emph{外延性（extensionality）:}
    %
    \begin{equation*}
      \fall{x, y : V} x \subseteq y \land y \subseteq x \Leftrightarrow x = y.
    \end{equation*}
    %
     \item \emph{空集（empty set）:} 对所有 $x:V$，均有 $\neg (x\in \emptyset)$。
    %
    \item \emph{无序对（pairing）:} 对所有 $u, v:V$，类 $\{u, v\} \defeq \setof{ x | x = u \vee x = v}$ 是一个 $V$-集。
      %
    \item \emph{无穷（infinity）:}\index{axiom!of infinity} 存在一个 $v:V$，满足 $\emptyset\in v$，且 $x\in v$ 蕴含 $x\cup \{x\}\in v$。
    %
  \item \emph{并集（union）:} 对所有 $v:V$，类 $\cup v\defeq \setof{ x | \exis{u:V} x \in u \in v}$ 是一个 $V$-集。
    %
    \item \emph{函数集（function set）:} 对所有 $u, v:V$，类 $v^u \defeq \setof{ x | x : u\to v}$ 是一个 $V$-集。%
      \footnote{此处 $x:u\to v$ 意指 $x$ 是一个适当的、由有序对所组成的集合，这是集合论中编码函数的惯常方式。}
    %
   \item \emph{$\in$-归纳（$\in$-induction）:} 如果 $C : V \to \prop$ 是一个类，使得只要对所有 $x\in a$ 都有 $C(x)$ 成立，则 $C(a)$ 便成立，那么对所有 $v:V$，均有 $C(v)$。
   %
     \item \emph{替换（replacement）:}\index{axiom!of replacement} 给定任意 $r : V \to V$ 和 $x : V$，类
       %
       \begin{equation*}
         \setof{ y | \exis{z : V} z \in x \land y = r(z)}
       \end{equation*}
       %
       是一个 $V$-集。
  %
   \item \emph{分离（separation）:}\index{axiom!of separation} 给定任意 $a : V$ 和 $\UU$-小的 $C : V \to \propU$，类
     %
     \begin{equation*}
       \setof{ x | x \in a \land C(x)}
     \end{equation*}
     %
     是一个 $V$-集。
  \end{enumerate}
\end{thm}

\begin{proof}[证明概要]
  \mbox{}
  %
  \begin{enumerate}
  \item 外延性（extensionality）：如果 $\vset(A,f) \subseteq \vset(B, g)$，那么对每个 $a : A$ 都有 $f(a) \in \vset(B, g)$，因此对每个 $a : A$ 仅有 $b:B$ 使得
    $f(a) = g(b)$。由假设 $\vset(B, g) \subseteq \vset(A, f)$ 可得~\eqref{eq:V-path} 的另一半，因此 $\vset(A,f) = \vset(B,g)$。

  \item 空集（empty set）：假设 $x\in \emptyset = \vset(\emptyt,\rec\emptyt(V))$。那么 $\exis{a:\emptyt}x=\, \rec\emptyt(V,a)$，这是荒谬的。

  \item 无序对（pairing）：给定 $u$ 和 $v$，令 $w=\vset(\bool,\rec\bool(V,u,v))$。
    \index{pair!unordered}

  \item 无穷（infinity）：取 $w = \vset(\nat,I)$，其中 $I: \nat \to V$ 由递归式 $I(0) \defeq \emptyset$ 及 $I(n+1) \defeq I(n)\cup \{I(n)\}$ 给出。

  \item 并集（union）：任取 $v:V$ 及其任意一个表示 $f :A\to V$ 满足 $v=\vset(A,f)$。然后令 $\tilde{A} \defeq \sm{a:A}[fa]$，其中 $m_{fa} : [fa] \mono V$ 是来自 \cref{def:TypeOfElements} 的成员类型。显然 $\tilde{A}$ 是 $\UU$-小的，并且我们有 $\cup v \defeq \vset(\tilde{A}, \lam{x} m_{f(\proj1(x))}(\proj2(x)))$。

  \item 函数集（function set）：给定 $u, v:V$，取成员类型 $[u] \mono V$ 和 $[v] \mono V$，以及函数类型 $[u]\to [v]$。我们想定义一个映射
  \[
 r: ([u]\to [v])\ \longrightarrow\ V
  \]
   使其满足“$r(f) = \setof{ \pairr{x, f(x)} | x : [u] }$”；但为使该式有意义，我们必须先定义有序对 $\pairr{x, y}$，然后取映射 $r': x \mapsto \pairr{x, f(x)}$，接着便可令 $r(f)\defeq \vset([u], r')$。我们可按惯常的方式，用无序对来定义有序对。

  \item $\in$-归纳（$\in$-induction）：令 $C : V \to \prop$ 为一个类，使得只要对所有 $x\in a$ 有 $C(x)$ 成立，则 $C(a)$ 成立；并任取 $v=\vset(B,g)$。为了通过归纳证明 $C(v)$，假设对所有 $b:B$ 有 $C(g(b))$。对任意 $x\in v$，仅有某个 $b:B$ 使得 $x = g(b)$，因此 $C(x)$ 成立。从而 $C(v)$ 得证。

  \item 替换（replacement）：用 $C$ 表示所讨论的类。
    断言“$C$ 是一个 $V$-集”是一个命题，因此我们可以如下进行归纳证明。
    假设 $x$ 是 $\vset(A, f)$，我们断言 $w \defeq \vset(A, r \circ f)$ 即为我们所寻找的集合。
    若 $C(y)$，则仅有 $z : V$ 和 $a : A$ 使得 $z = f(a)$ 且 $y = r(z)$，因此 $y \in w$。
    反之，若 $y \in w$，则仅有 $a : A$ 使得 $y = r(f(a))$，所以若取 $z \defeq f(a)$，便知 $C(y)$ 成立。

  \item 我们称一个类 $C: V\to\prop$ 是 \define{可分的}
    \indexdef{class!separable}%
    \indexdef{separable class}%
    如果对任意 $a:V$，类
  %
  \symlabel{class-intersection}
  \begin{equation*}
    a \cap C \defeq\setof{x | x\in a \wedge C(x)}
  \end{equation*}
  %
  是一个 $V$-集。
我们需要证明任何 $\UU$-小的 $C: V \to \propU$ 都是可分的。实际上，给定 $a=\vset(A,f)$，令 $A' = \sm{x:A}C(fx)$，并取 $f' = f\circ i$，其中 $i : A' \to A$ 是显然的包含映射。然后我们可以取 $a' = \vset(A',f')$，并且 $x\in a\wedge C(x) \Leftrightarrow x\in a'$ 如所断言。我们需要 $C$ 取值于 $\UU$ 的假设，以便 $A' = \sm{x:A}C(fx)$ 能落在 $\UU$ 中。\qedhere
\end{enumerate}
\end{proof}

另外，拥有一个关于可分性的严格句法判据也很方便，这样从类的表达式就能直接看出它产生的是一个 $V$-集。一个熟悉的此类条件就是“$\Delta_0$”，它意味着表达式仅由相等关系 $x=_V y$ 和隶属关系 $x\in y$ 经由纯命题联结词 $\neg$、$\land$、$\lor$、$\Rightarrow$ 以及关于特定集合的量词（即形如 $\exists(x\in a)$ 和 $\forall(y\in b)$ 的量词，这些称为\define{有界的（bounded）}量词\index{bounded!quantifier}\index{quantifier!bounded}）构造而成。\indexdef{separation!.Delta0@$\Delta_0$}%

\begin{cor}\label{cor:Delta0sep}
如果类 $C: V \to \prop$ 是上述意义上的 $\Delta_0$ 类，那么它是可分的。
\end{cor}


\index{axiom!of $\Delta_0$-separation}%

\begin{proof}
回顾一下，$x = y$ 有一个 $\UU$-小的调整 $x \bisim y$。由于 $x\in y$ 是通过 $x=y$ 定义的，我们同样有隶属关系的一个 $\UU$-小的调整
%
\symlabel{resized-membership}
\begin{equation*}
  x\bin\vset(A,f) \defeq \exis{a:A} x \bisim f(a).
\end{equation*}
%
现在，令 $\Phi$ 为 $C$ 的一个 $\Delta_0$ 表达式，使得作为类有 $\Phi = C$（严格来说，我们应当区分表达式与其含义，但这里我们模糊这一区别）。令 $\widetilde{\Phi}$ 为将其中所有 $=$ 和 $\in$ 的出现替换为调整后的对应物 $\bisim$ 和 $\bin$ 所得的结果。显然 $\widetilde{\Phi}$ 也表示 $C$，即对所有 $x:V$ 有 $\widetilde{\Phi}(x) \Leftrightarrow C(x)$，因而由泛等性得 $\widetilde{\Phi}=C$。现在只需证明 $\widetilde{\Phi}$ 是 $\UU$-小的，那么根据定理，它便是可分的。

我们通过对表达式的构造进行归纳来证明 $\widetilde{\Phi}$ 是 $\UU$-小的。基础情形是 $x \bisim y$ 和 $x\bin y$，它们已被调整到 $\UU$ 中。$\UU$ 显然在单纯命题运算（以及 $(-1)$-截断）下封闭，于是只剩下有界量词 $\exists(x\in a)$ 和 $\forall(y\in b)$ 需要检查。根据定义，
\begin{align*}
\exists(x\in a) P(x) &\defeq \Brck {\sm{x:V}(x\bin a \land P(x))},\\
\forall(y\in b) P(x) &\defeq  \prd{x:V}(x\bin a \to P(x)).
\end{align*}
考虑 $\brck {\sm{x:V}(x\bin a \land P(x))}$。根据归纳假设 $P(x)$ 是 $\UU$-小的，从而主体 $(x\bin a \land P(x))$ 是 $\UU$-小的，但对 $V$ 的量化不一定保持在 $\UU$ 内。然而，在当前情形下，我们可以将其替换为对 $a$ 的成员类型 $[a]\mono V$ 的量化，且不难证明
\begin{equation*}
  \sm{x:V}(x\bin a \land P(x)) = \sm{x:[a]} P(x).
\end{equation*}
由于 $[a]$ 和 $P(x)$ 都在 $\UU$ 中，右端确实仍在 $\UU$ 中。$\prd{x:V}(x\bin a \to P(x))$ 的情形类似，利用 $\prd{x:V}(x\bin a \to P(x)) = \prd{x:[a]}P(x)$ 即可。
\end{proof}

我们已经证明，在带有一个宇宙 $\UU$ 的类型论中，累积层次 $V$ 是一个包含许多标准公理的"构造性集合论"
\index{constructive!set theory}%
的模型。
然而，据我们所知，它缺少 \emph{强收集}
\index{axiom!strong collection}%
\index{collection!strong}%
\index{strong!collection}%
与 \emph{子集收集}
\index{axiom!subset collection}%
\index{collection!subset}%
\index{subset!collection}%
公理，而这两条公理包含在 \CZF{}~\cite{AczelCZF} 中。
在这种集合论到类型论的通常解释中，这两条公理是类似集囿的相等定义的推论；而在其他构造出的集合论模型中，强收集可能因其他原因而成立。
我们不知道这两条公理是否在我们的模型 $(V,\in)$ 中成立，但似乎不太可能。
由于 $V$ 是系统\emph{内}的高阶归纳类型，而非\emph{外部}构造，它在某些方面与以往的解释不同并不令人意外。

最后，考虑将集合的选择公理以 $\choice{}$ 的形式（来自上文的 \cref{subsec:emacinsets}）添加到我们的类型论中。根据 \cref{thm:1surj_to_surj_to_pem}，可得 $\LEM{}$ 也成立，进而由 \cref{thm:settopos} 知 $\set$ 是一个带有子对象分类子 $\bool$ 的意象\index{topos}。此时，有 $\prop = \bool:\UU$，因而\emph{所有类都是可分的}。
于是我们证明了：

\begin{lem}\label{lem:fullsep}
  在带有 $\choice{}$ 的类型论中，$V$ 满足 \define{（完全）分离}
  \indexdef{separation!full}%
  律：给定\emph{任意}类 $C : V \to \prop$ 和 $a : V$，类 $a \cap C$ 是一个 $V$-集。
\end{lem}

\begin{thm}\label{thm:zfc}
在带选择公理 $\choice{}$ 和宇宙 $\UU$ 的类型论中，累积层次 $V$ 是带选择公理的 Zermelo--Fraenkel\index{set theory!Zermelo--Fraenkel} 集合论（ZFC）的一个模型。
\end{thm}

\begin{proof}
我们已经具备了 \cref{thm:VisCST} 所列的所有公理，以及完全分离律，因此只需证明对每个 $a:V$ 都存在幂集\index{power set} $\power a:V$。
但由于我们有排中律 $\LEM{}$，这些幂集即为函数类型 $\power a = (a\to\bool)$。
于是 $V$ 便成为 Zermelo--Fraenkel 集合论（ZF）的一个模型。
从 $\choice{}$ 验证集合论的选择公理，我们留作一个简单的练习。
\end{proof}

\index{bargaining|)}%

%%%%%%%%%%%%%%%%%%%%%%%%%%%%%%%%%%%%%%%%%%%%%%%%%%%%%%%%%%%%%%%%%%%%%%
\sectionNotes

人们对集合范畴所期望的基本性质，可以追溯到初等意象理论的早期。
在 \cref{subsec:emacinsets} 中提及的\emph{集合范畴的初等理论（Elementary theory of the category of sets）}是由 Lawvere\index{Lawvere} 在 \cite{lawvere:etcs-long} 中引入的，作为对集合论的一种范畴论公理化。
\index{Elementary Theory of the Category of Sets}%
$\Pi W$-预意象（$\Pi W$-pretopos）的概念，被视为初等意象的一个直谓版本，是在~\cite{MoerdijkPalmgren2002} 中提出的；亦可参阅~\cite{palmgren:cetcs}。

\cref{sec:piw-pretopos} 中对集合范畴的处理大体遵循了~\cite{RijkeSpitters}。
满态射是满射这一事实（\cref{epis-surj}）在经典数学中广为人知，但要\emph{直谓地}证明它却并不像看起来那么平凡。
\index{数学!直谓的}%
\cite{Mines/R/R:1988} 中的证明用到了幂集运算（这一运算本身是非直谓的），不过，该证明也可以看作如下较弱命题的一个直谓证明：若一个映射在宇宙 $\UU_i$ 中是满态射，则它在下一个宇宙 $\UU_{i+1}$ 中是满射。
Wilander~\cite{Wilander2010} 给出了关于集囿（setoids）的一个直谓证明。
我们的证明与 Wilander 的类似，但通过使用推出和泛等性避开了集囿。

\cref{thm:1surj_to_surj_to_pem} 中从 $\choice{}$ 推出 $\LEM{}$ 的蕴含关系，是 Diaconescu~\cite{Diaconescu} 在意象理论中证明的一个定理在同伦类型论中的改编；Bishop~\cite[Problem~2]{Bishop1967} 早已将其作为一个问题提出。

关于序数的直觉主义理论，参见~\cite{taylor:ordinals,Taylor99} 以及 \cite{JoyalMoerdijk1995}。
在类型论中通过归纳原理（包括可达性\index{accessibility}这一归纳谓词）定义良基性的研究，见于~\cite{Huet80,Paulson86,Nordstrom88}，尽管这一思想可以追溯到 Gentzen 对算术一致性\index{一致性!算术的}的证明~\cite{Gentzen36}。

代数集合论（algebraic set theory）的思想启发了我们在 \cref{sec:cumulative-hierarchy} 中对累积层次的发展，这一思想归功于~\cite{JoyalMoerdijk1995}，但它源自更早的工作~\cite{AczelCZF}。
\index{algebraic set theory}%
\index{集合论!代数的}%

%%%%%%%%%%%%%%%%%%%%%%%%%%%%%%%%%%%%%%%%%%%%%%%%%%%%%%%%%%%%%%%%%%%%%%
\sectionExercises

\begin{ex}\label{ex:utype-ct}
  仿照 $\uset$ 的模式，我们希望构造一个由所有类型（在给定宇宙 $\UU$ 中）以及它们之间的映射组成的范畴 $\utype$。
  然而，为了使它成为 \cref{sec:cats} 意义上的范畴，我们必须先声明 $\hom(X,Y) \defeq \pizero{X\to Y}$，其上的复合由截断归纳从普通复合 $(Y\to Z) \to (X\to Y) \to (X\to Z)$ 来定义。
  这在 \cref{ct:hoprecat} 中被定义为\emph{类型的同伦预范畴（homotopy precategory of types）}。
  但它仍然不是一个范畴，而只是一个预范畴（其对象类型 $\UU$ 甚至连 $0$-类型都不是）。
  经过 Rezk 完备化
  \index{completion!Rezk}%
  （参见 \cref{ct:hocat}），它成为一个范畴，并且根据 \cref{ct:ex:hocat}，其对象类型可以等同于 $\trunc1\type$。
  证明所得到的范畴 $\utype$ 与 $\uset$ 不同，它不是一个预意象。
\end{ex}

\begin{ex}\label{ex:surjections-have-sections-impl-ac}
  证明：如果在范畴 $\uset$ 中每个满射都有一个截面，那么选择公理成立。
\end{ex}

\begin{ex}\label{ex:well-pointed}
  证明：在 $\LEM{}$ 下，范畴 $\uset$ 是良点的，
  \indexdef{category!well-pointed}%
  即以下陈述成立：对于任意 $f, g : A\to B$，如果 $f \neq g$，则存在函数 $a : 1\to A$ 使得 $f(a) \neq g(a)$。
  证明切片范畴
  \index{category!slice}%
  $\uset/\bool$（由函数 $A\to \bool$ 和交换三角形构成）不具备此性质。
  （提示：$\uset/\bool$ 中的终对象是恒等函数 $\bool \to \bool$，因此在这个范畴中，存在没有元素 $1\to X$ 的对象 $X$。）
\end{ex}

\begin{ex}\label{ex:add-ordinals}
  \index{addition!of ordinal numbers}%
  证明：如果 $(A,<_A)$ 和 $(B,<_B)$ 是良基的、外延的，或者是序数，那么 $A+B$ 同样如此，其中 $<$ 定义为
  \begin{align*}
    (a<a') &\defeq (a<_A a') & \text{for }& a,a':A\\
    (b<b') &\defeq (b<_B b') & \text{for }& b,b':B\\
    (a<b) &\defeq \unit      & \text{for }& (a:A),(b:B)\\
    (b<a) &\defeq \emptyt    & \text{for }& (a:A),(b:B).
  \end{align*}
\end{ex}


% \begin{proof}
%   We first prove by induction on $<_A$ that every element of $A$ is accessible in $A+B$.
%   This is easy since the only elements less than $a:A$ in $A+B$ are also in $A$.
%   We then prove by induction on $<_B$ that every element of $B$ is accessible in $A+B$.
%   This is easy since we have already proven that every element of $A$ is accessible.
% \end{proof}

\begin{ex}\label{ex:multiply-ordinals}
  \index{multiplication!of ordinal numbers}%
  证明：如果 $(A,<_A)$ 和 $(B,<_B)$ 是良基的、外延的，或者是序数，那么 $A\times B$ 同样如此，其中 $<$ 定义为
  \[ ((a,b) <(a',b')) \defeq (a<_A a') \vee ((a=a') \wedge (b<_B b')). \]
\end{ex}


% \begin{proof}
%   We prove by induction on $<_A$ that for every $a:A$, every element of the form $(a,b)$ is accessible in $A\times B$.
%   The inductive hypothesis is that for all $a'<_A a$, every pair $(a',b)$ is accessible.
%   Inside this induction, we prove by induction on $<_B$ that for every $b:B$, the element $(a,b)$ is accessible.
%   The nested inductive hypothesis is that for every $b'<_B b$, the element $(a,b')$ is accessible.
%   But now, if $(a',b')< (a,b)$, then either $a<_A a'$ in which case $(a',b')$ is accessible by the first inductive hypothesis, or $a=a'$ and $b'<_B b$, in which case $(a,b')$ is accessible by the second inductive hypothesis.
%   Thus, by definition of accessibility, $(a,b)$ is accessible.
%   This completes both inductions.
% \end{proof}

\begin{ex}\label{ex:algebraic-ordinals}
  在序数上定义通常的代数运算，并证明它们满足通常的性质。
\end{ex}

\begin{ex}\label{ex:prop-ord}
  注意 $\bool$ 是一个序数，它带有一个显然的 $<$ 关系，其中仅 $\bfalse<\btrue$ 成立。
  \begin{enumerate}
  \item 在 $\prop$ 上定义一个 $<$ 关系，使其成为一个序数。
  \item 证明 $\id[\ord]\bool\prop$ 当且仅当 \LEM{} 成立。
  \end{enumerate}
\end{ex}

\begin{ex}\label{ex:ninf-ord}
  回顾，当把 $\nat$ 视为序数时，我们称其为 $\omega$；因此我们也有序数 $\omega+1$。
  另一方面，我们定义
  \[ \nat_\infty \defeq \setof{a:\nat\to\bool | \fall{n:\nat} (a_n \le a_{\suc(n)}) } \]
  其中 $\le$ 表示 $\bool$ 上显然的偏序关系，满足 $\bfalse\le\btrue$。
  \begin{enumerate}
  \item 在 $\nat_\infty$ 上定义一个 $<$ 关系，使其成为一个序数。
  \item 证明 $\id[\ord]{\omega+1}{\nat_\infty}$ 当且仅当有限全知原理~\eqref{eq:lpo} 成立。%
    \index{limited principle of omniscience}%
  \end{enumerate}
\end{ex}

\begin{ex}\label{ex:well-founded-extensional-simulation}
  证明：若 $(A,<)$ 是良基且外延的，并且 $A:\UU$，则存在一个模拟 $A\to V$，其中 $(V,\in)$ 是从宇宙~\UU 构建的累积层次（见 \cref{sec:cumulative-hierarchy}）。
\end{ex}

\begin{ex}\label{ex:choice-function}
  证明 \cref{thm:wop}\ref{item:wop1} 等价于选择公理~\eqref{eq:ac}。
\end{ex}

\begin{ex}\label{ex:cumhierhit}
  给定类型 $A$ 和 $B$，定义一个\define{双完全关系}
  \indexsee{bitotal relation}{relation, bitotal}%
  \indexdef{relation!bitotal}%
  为 $R:A\to B\to \prop$，使得
  \[ \Big(\fall{a:A}\exis{b:B} R(a,b) \Big) \land \Big(\fall{b:B}\exis{a:A} R(a,b) \Big). \]
  对于这样的 $A$、$B$、$R$，令 $A\sqcup^R B$ 为由以下构造子生成的高阶归纳类型：
  \begin{itemize}
  \item $i:A\to A\sqcup^R B$，
  \item $j:B\to A\sqcup^R B$，
  \item 对每一对满足 $R(a,b)$ 的 $a:A$ 和 $b:B$，有一条道路 $i(a)=j(b)$。
  \end{itemize}
  证明累积层次 $V$ 可以用以下更直接的构造子列表来定义，并且所得的归纳原理就是 \cref{sec:cumulative-hierarchy} 中给出的那个：
  \begin{itemize}
  \item 对每个 $A : \UU$ 和 $f : A \to V$，有一个元素 $\vset(A, f) : V$。
  \item 对任意 $A,B:\UU$ 和双完全关系
    \index{relation!bitotal}%
    $R:A\to B\to \prop$，以及任意映射 $h:A\sqcup^R B \to V$，存在一条道路 $\id{\vset(A,h\circ i)}{\vset(B,h\circ j)}$。
  \item 0-截断构造子。
  \end{itemize}
\end{ex}

\begin{ex}\label{ex:strong-collection}
  在 \CZF 中，\define{强收集公理}
  \indexdef{axiom!strong collection}%
  \indexdef{collection!strong}%
  \indexdef{strong!collection}%
  的形式为：
   \begin{multline*}
   \Parens{\fall{x\in v}\exis{y} R(x,y)} \Rightarrow \\
   \exis{w}\big[\big(\fall{x\in v}\exis{y\in w}R(x,y)\big)\land \big(\fall{y\in w}\exis{x\in v}R(x,y) \big)\big]
   \end{multline*}
   它在累积层次 $V$ 中成立吗？（我们不知道这个问题的答案。）
\end{ex}

\begin{ex}\label{ex:choice-cumulative-hierarchy-choice}
验证：若假定 $\choice{}$，则累积层次 $V$ 满足通常集合论意义上的选择公理，该公理可以表述为：
  \[
   \fall{x:V} \Parens{(\fall{y\in x}\exis{z:V} z\in y) \Rightarrow  \exis{c\in(\cup x)^x}\fall{y\in x} c(y)\in y}
   \]
\end{ex}

\begin{ex}\label{ex:plump-ordinals}
  假定命题大小重整，证明存在一个单纯的谓词 $\mathsf{isPlump}:\ord\to\prop$，使得对于任意 $A:\ord$ 有
  \begin{multline*}\label{eq:plump}
    \mathsf{isPlump}(A) = \Parens{\fall{B<A} \mathsf{isPlump}(B)} \wedge\narrowbreak
    \Parens{\fall{C,B:\ord} C\le B < A \wedge \mathsf{isPlump}(C) \Rightarrow C < A}.
  \end{multline*}
  注意，$\mathsf{isPlump}$ 不能简单地通过关于 \ord 的良基归纳来定义；你必须使用一个不同的良基关系。
  我们说一个序数 $A$ 是\define{丰满的}~\cite{taylor:ordinals,Taylor99}，如果 $\mathsf{isPlump}(A)$。
  \index{plump!ordinal}\index{ordinal!plump}%
\end{ex}

\begin{ex}\label{ex:not-plump}
  证明 \LEM{} 等价于"所有序数都是丰满的"这一陈述。
\end{ex}

\begin{ex}\label{ex:plump-successor}
  定义序数 $A$ 的\define{丰满后继}\index{plump!successor}\indexdef{successor!plump}为
  \[ t(A) \defeq \setof{ B:\ord | (B\le A) \wedge \mathsf{isPlump}(B) } \]
  \begin{enumerate}
  \item 根据定义，$t(A)$ 属于下一个更高的宇宙。证明：在假定命题大小重整的条件下，它等于一个与 $A$ 在同一宇宙中的序数。
  \item 再次假定命题大小重整，证明如果 $A$ 是丰满的（\cref{ex:plump-ordinals}），那么 $t(A)$ 也是丰满的。
  \end{enumerate}
\end{ex}

\begin{ex}\label{ex:ZF-algebras}
  相对于宇宙 $\UU_i$ 的一个\define{ZF-代数}~\cite{JoyalMoerdijk1995}
  \index{ZF-algebra}%
  是一个偏序集（见 \cref{ct:orders}）$V:\UU_{i+1}$，它具有所有以 $\UU_i$ 中类型为指标的上确界，并且配备了一个"后继"函数 $s:V\to V$（不一定在任何意义上保持 $\le$）。
  \begin{enumerate}
  \item 证明累积层次 $(V_{\UU_i},\subseteq,s)$ 是初始 ZF-代数，其中 $s(x)$ 是单点集 $\setof{x}$。
  \item 证明 $(\ord_{\UU_i},\le,s)$ 是满足"对所有 $x$ 有 $x\le s(x)$"这一性质的初始 ZF-代数，其中 $s(A)=A+\unit$ 是来自 \cref{thm:ordsucc} 的后继\index{successor!of an ordinal}。
  \item 假定命题大小重整，证明 $\Parens{\setof{A:\ord_{\UU_i} | \mathsf{isPlump}(A) },\le,t}$ 是满足"对所有 $x,y$ 有 $(x\le y) \Rightarrow (t(x)\le t(y))$"这一性质的初始 ZF-代数，其中 $t$ 是来自 \cref{ex:plump-successor} 的丰满后继。
  \end{enumerate}
\end{ex}

\begin{ex}\label{ex:monos-are-split-monos-iff-LEM-holds}
  对于范畴 $A$，态射 $f: \hom_A(a,b)$ 被称为是一个\define{可分裂单态射}，如果存在一个态射 $g: \hom_A(b,a)$ 使得 $\id{g \circ f}{1_a}$。（这样的 $g$ 称为 $f$ 的一个\define{收缩}。）
  证明以下陈述逻辑等价：
  \begin{enumerate}
  \item $\LEM{}$。
  \item 对所有集合 $A$ 和 $B$，如果 $A$ 是有实例的，那么对于 $\uset$ 中的每一个单态射 $f: A \to B$，$f$ 在 $\uset$ 中也是可分裂单态射。
  \end{enumerate}
\end{ex}

\index{set|)}%

% Local Variables:
% TeX-master: "hott-online"
% End: